\chapter{约束动力学与解析微分：从 DAE 到可微接触与 MPC 梯度}
\label{ch:constrained-dynamics}

% ════════════════════════════════════════════════════════════════
%  控制部 C2「刚体动力学」子部末章。承 J(空间向量/RNEA/ABA)、K(Lagrange/Hamilton/辛)，
%  启 N(接触)、O/P(操作空间/WBC)、H(DDP)、G(MPC)。
%  两源融合：Tutorial 60_约束动力学 + 50_动力学解析微分（范围 A：理论主体全保留，纯工程跳）。
%  记号：右扰动 J_r 主、se(3) ξ=[rho;φ]、KKT 约束力 +J^⊤lambda、Delassus G=JM^{-1}J^⊤、Λ=G^{-1}。
% ════════════════════════════════════════════════════════════════

\begin{practice}[前置自测]
开始之前先自测。若已能流畅作答，可速读本章的指数阶梯表、KKT 速查与解析导数递推；若卡住，正是本章要打通的。
\begin{enumerate}
  \item 写出固定基机器人动力学 $\mathbf M(\mathbf q)\ddot{\mathbf q}+\mathbf C(\mathbf q,\dot{\mathbf q})\dot{\mathbf q}+\mathbf g(\mathbf q)=\boldsymbol\tau$ 中每一项的物理含义与维度（\cref{ch:geometric-mechanics} 已给出）。
  \item 什么是 Frobenius 定理？分布可积的等价条件是什么（\cref{ch:lie}）？
  \item 解释 KKT 互补松弛 $\lambda_i g_i=0$ 的物理含义。
  \item 什么是微分代数方程（DAE）？它与 ODE 的本质区别是什么？
  \item DDP/iLQR 的 Riccati 反传需要哪两个动力学 Jacobian？对 $36$ 自由度系统，用中心差分算它们要调用多少次 RNEA（\cref{ch:spatial-dynamics}）？
\end{enumerate}
\end{practice}

\section*{本章目标}
学完本章，你应当能够：
\begin{enumerate}
  \item \textbf{形式化}任意约束系统为 DAE 并识别其微分指数；用 Baumgarte 稳定化消除数值漂移；
  \item \textbf{推导}从 d'Alembert 原理到完整 KKT 鞍点系统的每一步，理解 Lagrange 乘子即约束力；
  \item \textbf{建模}闭链机构（Stewart/Delta/Cassie 四杆膝）的环路闭合约束并用 Schur 补求解；
  \item \textbf{理解}接触动力学的 LCP/NCP 表述、Painlevé 悖论与四条可微化路线；
  \item \textbf{掌握}约束动力学的统一视角——Gauss 最小约束原理、Udwadia--Kalaba 显式解、IFT-on-KKT；
  \item \textbf{推导} $\partial\mathrm{RNEA}/\partial\mathbf q$ 的 $O(Nd)$ 解析递推（placement 微分恒等式）与 $\partial\mathrm{ABA}$ via ID--FD 互逆；
  \item \textbf{区分}一阶/二阶解析导数、惯性辨识 regressor，并把它们焊接到 DDP/MPC 的梯度链上。
\end{enumerate}

\subsection*{承上与本章定位}
\cref{ch:spatial-dynamics} 用空间向量代数把固定基机器人的逆动力学（RNEA）与正动力学（ABA）压到 $O(N)$，\cref{ch:geometric-mechanics} 从 Lagrange/Hamilton 与 $\SE(3)$ 几何力学给出了 $\mathbf M,\mathbf C,\mathbf g$ 的结构与辛/守恒律。但那里的机器人都\emph{在真空中自由漂浮}——它没有关节极限、没有足端踩地、没有闭链联动、没有抓握形成的闭合回路。现实机器人处处是约束，而\textbf{约束不是给方程加一个条件，它改变了方程的数学结构}：从 ODE 变为 DAE，从光滑流变为非光滑混合系统，从可逆线性系统变为互补问题。本章正是\cref{ch:control}「刚体动力学」子部的收口章，把这条「约束改变结构 $\to$ KKT 求解 $\to$ 解析微分喂最优控制」的主线一次打通。它在全书的独特之处是\emph{数学跨度最大}：从 18 世纪的分析力学（d'Alembert 1743、Lagrange 1788、Gauss 1829）一直走到 21 世纪的可微仿真（Pinocchio~3、Dojo、Nimble），用同一套语言串起来。

\subsection*{本章主线}
全章分四部分，循一条递进的叙事：

\begin{center}
\small
\begin{tikzpicture}[
  node distance=3.6mm and 6mm, >=Stealth,
  blk/.style={draw, rounded corners, align=center, font=\small, inner sep=3pt, fill=main!6, draw=main!60!black},
  grp/.style={draw, rounded corners, align=center, font=\small, inner sep=3pt, fill=second!6, draw=second!60!black},
  prp/.style={draw, rounded corners, align=center, font=\small, inner sep=3pt, fill=third!8, draw=third!60!black},
  ln/.style={->, thick, gray!60!black}]
\node[blk] (cls) {约束分类\\Frobenius/Signorini};
\node[blk, right=of cls] (dae) {DAE/指数\\drift-off};
\node[blk, right=of dae] (kkt) {KKT 鞍点\\Schur/Delassus};
\node[grp, below=7mm of cls] (baum) {Baumgarte\\临界阻尼};
\node[grp, right=of baum] (loop) {闭链\\cut-joint};
\node[grp, right=of loop] (cont) {接触\\LCP/摩擦锥};
\node[prp, below=7mm of baum] (gauss) {Gauss/UK\\IFT-on-KKT};
\node[prp, right=of gauss] (diff) {可微接触\\四路线};
\node[prp, right=of diff] (ad) {解析微分\\$\partial$RNEA/$\partial$ABA};
\draw[ln] (cls)--(dae); \draw[ln] (dae)--(kkt);
\draw[ln] (kkt.south)-- ++(0,-3.5mm) -| (baum.north);
\draw[ln] (baum)--(loop); \draw[ln] (loop)--(cont);
\draw[ln] (cont.south)-- ++(0,-3.5mm) -| (gauss.north);
\draw[ln] (gauss)--(diff); \draw[ln] (diff)--(ad);
\draw[ln] (ad.east) -- ++(6mm,0) node[right, font=\scriptsize, align=left] {喂\\DDP/MPC};
\end{tikzpicture}
\end{center}

\begin{note}
\textbf{本章记号与约定.}\;
广义坐标 $\mathbf q\in\mathcal Q$，速度维度 $n_v$（浮基时 $n_v=n+6$、$n_q\neq n_v$）。约束 Jacobian $\mathbf J=\partial\boldsymbol\varphi/\partial\mathbf q\in\mathbb R^{m\times n_v}$，接触 Jacobian 记 $\mathbf J_c$。延续 \cref{ch:spatial-dynamics}，空间量（$\mathbf v_i,\mathbf a_i,\mathbf f_i\in\mathbb R^6$、运动子空间 $\mathbf S_i$、Plücker 变换 ${}^i\mathbf X_{\lambda(i)}$、空间叉乘 $\times/\times^*=\mathrm{ad}/\mathrm{ad}^*$）一律用 Featherstone 的 $[\boldsymbol\omega;\mathbf v]$ 序；而\textbf{对 $\mathbf q$ 的求导一律取本书右扰动} $\mathbf R\,\mathrm{Exp}(\delta\boldsymbol\phi)$（右雅可比 $\mathbf J_r$，与 \cref{ch:lie} 一致），二者由置换 $\mathbf P=\begin{psmallmatrix}\mathbf 0&\mathbf I\\\mathbf I&\mathbf 0\end{psmallmatrix}$ 关联（$\se(3)$ 序 $\boldsymbol\xi=[\boldsymbol\rho;\boldsymbol\phi]$）。\textbf{约束力恒为 Jacobian 的转置乘乘子 $\mathbf J^\top\boldsymbol\lambda$}（缺转置则力方向错），KKT 用 $\begin{psmallmatrix}\mathbf M&-\mathbf J^\top\\\mathbf J&\mathbf 0\end{psmallmatrix}$。\textbf{Delassus 矩阵 $\mathbf G:=\mathbf J\mathbf M^{-1}\mathbf J^\top$、操作空间惯量 $\boldsymbol\Lambda:=\mathbf G^{-1}$}。摩擦系数记 $\mu_f$，proximal 正则化记 $\mu$（二者不同，务必区分）。
\end{note}

% ════════════════════════════════════════════════════════════════
\section{约束如何改变方程结构}
\label{sec:cd-motivation}

\paragraph{动机：现实约束无处不在。}
固定基机械臂之外，几乎所有机器人都带约束：足端非穿透与摩擦（腿足），抓握物体形成的闭合回路（操作），并联机构与四杆膝的环路闭合（Stewart/Delta/Cassie），以及无处不在的关节限位。这些「限制」绝非同一类——把它们一视同仁地用同一种方法处理会立刻翻车：对单边接触强行当双边约束，足离地时会冒出非物理的「吸力」把脚拉回地面；对不可积的轮子约束试图积出位置层约束，根本做不到；对闭链不分清约束维度，KKT 矩阵尺寸就错了。

\begin{insight}[约束改变的是数学结构]
约束不是「给方程加一个条件」那么轻巧。它把系统从 ODE 推到 DAE，把光滑流推到非光滑混合系统，把可逆线性求解推到互补问题。理解约束的\emph{分类}，就是先认清你面对的数学问题属于哪一类——这决定了后面用 KKT 还是 LCP、用 Baumgarte 还是投影、用 dense 分解还是 proximal。
\end{insight}

本章的全景已在主线图中：先做约束分类学（\cref{sec:cd-classify}），把一阶约束统一到 Pfaffian 形式并用 Frobenius/Signorini 判型；再形式化为 DAE、揭示微分指数与 drift-off（\cref{sec:cd-dae}）；接着从 d'Alembert 推出 KKT 鞍点（\cref{sec:cd-kkt}）、给出 Baumgarte 稳定化及其替代（\cref{sec:cd-baumgarte}）；然后处理闭链（\cref{sec:cd-loop}）与接触（\cref{sec:cd-contact}、\cref{sec:cd-cwc}）；登上 Gauss/UK/IFT 的统一高地（\cref{sec:cd-gauss}）与可微接触四路线（\cref{sec:cd-diffcontact}）；最后转入动力学解析微分（\cref{sec:cd-deriv-need}--\cref{sec:cd-second}），把所有导数焊到 DDP/MPC 上。

% ════════════════════════════════════════════════════════════════
\section{约束分类学}
\label{sec:cd-classify}

\paragraph{物理直觉先行。}
想象一座山丘的等高线。\emph{完整约束}像「只能沿等高线走」——你被困在一张低维曲面上，能去的位置变少了。\emph{非完整约束}像「每一步只能朝北或朝东」——它不限制你最终能到哪些位置（仍可走遍整个平面），只限制你的\emph{瞬时速度方向}。把这条直觉形式化，就得到约束分类的两条主轴：可积性与单边性。

\subsection{Pfaffian 形式与 Frobenius 判据}
设广义坐标 $\mathbf q\in\mathcal Q\subset\mathbb R^n$。一切一阶约束统一写成 Pfaffian 形式
\begin{equation}
  \mathbf A(\mathbf q)\dot{\mathbf q}+\mathbf b(\mathbf q,t)=\mathbf 0,\qquad \mathbf A(\mathbf q)\in\mathbb R^{m\times n}.
  \label{eq:cd-pfaffian}
\end{equation}
若存在 $\boldsymbol\varphi(\mathbf q,t)$ 使 \cref{eq:cd-pfaffian} 等价于 $\dot{\boldsymbol\varphi}=\mathbf 0$（即 $\mathbf A\,\mathrm d\mathbf q+\mathbf b\,\mathrm dt$ 为恰当 1-形式），则约束\textbf{完整}（holonomic），可降维为位置层约束 $\boldsymbol\varphi(\mathbf q,t)=\mathbf 0$；否则\textbf{非完整}（nonholonomic）。

\begin{theorem}[可积判据 = Frobenius\cite{lynch2017modern}]\label{thm:cd-frobenius}
对分布 $\Delta=\ker(\mathbf A)$，
\[
  \Delta\ \text{可积}\iff\Delta\ \text{involutive}:\ \forall\,X,Y\in\Gamma(\Delta),\ [X,Y]\in\Gamma(\Delta),
\]
等价地 1-形式系统满足 $\mathrm d\omega_i\wedge\omega_1\wedge\cdots\wedge\omega_m=0\ (\forall i)$。
\end{theorem}
\noindent 此即 \cref{ch:lie} 微分几何的 Frobenius 定理在约束上的直接应用：可积 $\Leftrightarrow$ 约束分布封闭于 Lie 括号 $\Leftrightarrow$ 运动被困在子流形上。\cref{thm:cd-frobenius} 是本章约束分类的判据基石，值得把它\emph{证透}——尤其难方向“对合 $\Rightarrow$ 可积”是全部「非完整 vs 完整」之分的来源。下面给出完整证明，再把它落到单约束的一行可算检验上。

\begin{proof}[Frobenius 定理证明：对合 $\Leftrightarrow$ 可积]
记 $\Delta$ 为常秩 $k$ 的正则分布，$n=\dim\mathcal Q$，$\Gamma(\Delta)$ 为取值于 $\Delta$ 的光滑向量场。

\textbf{（易方向）可积 $\Rightarrow$ 对合。}\;
设 $\Delta$ 可积：过每点 $\mathbf q$ 有一张 $k$ 维积分流形 $\mathcal N$，满足 $T_{\mathbf p}\mathcal N=\Delta_{\mathbf p}$。取 $X,Y\in\Gamma(\Delta)$，二者处处切于叶 $\mathcal N$，其流 $\phi^X_t,\phi^Y_s$ 把 $\mathcal N$ 映到自身（始于叶、切于 $\Delta$ 的曲线不离开该叶）。于是 $X,Y$ 限制在 $\mathcal N$ 上是 $\mathcal N$ 的内蕴向量场，而李括号是子流形的\emph{内蕴}运算，故 $[X,Y]_{\mathbf p}\in T_{\mathbf p}\mathcal N=\Delta_{\mathbf p}$，即 $[X,Y]\in\Gamma(\Delta)$。等价的坐标视角：可积给出适配坐标 $(x^1,\dots,x^n)$ 使叶为 $\{x^{k+1}=c_{k+1},\dots,x^n=c_n\}$、$\Delta=\operatorname{span}\{\partial_{x^1},\dots,\partial_{x^k}\}$；此时 $X=\sum_{a\le k}X^a\partial_{x^a}$、$Y=\sum_{b\le k}Y^b\partial_{x^b}$，其括号第 $c\,(c>k)$ 分量 $[X,Y]^c=\sum_i(X^i\partial_i Y^c-Y^i\partial_i X^c)=0$（因 $X^c=Y^c\equiv0$），故 $[X,Y]\in\Delta$。

\textbf{（难方向）对合 $\Rightarrow$ 可积。}\;
核心是构造一组\emph{两两对易}的局部标架张成 $\Delta$，再同时拉直为坐标向量场。

\emph{第 1 步（化为阶梯标架）.}\, 在 $\mathbf q$ 附近取坐标 $(u^1,\dots,u^n)$。$\Delta$ 的任一标架 $\{X_1,\dots,X_k\}$ 给出 $k\times n$ 分量矩阵；因 $X_i$ 线性无关，必有 $k$ 列在 $\mathbf q$ 处构成可逆子块，重排坐标使前 $k$ 列可逆。左乘该 $k\times k$ 子块的逆（光滑、可逆的标架变换，不改变 $\Delta$），得规范化标架
\[
  Y_i=\partial_{u^i}+\sum_{a>k}b_i^{\,a}(\mathbf u)\,\partial_{u^a},\qquad i=1,\dots,k,
\]
即每个 $Y_i$ 的“前 $k$ 分量”恰为 $\delta_{i\cdot}$、其余分量沿 $\partial_{u^{k+1}},\dots,\partial_{u^n}$。

\emph{第 2 步（对合迫使对易）.}\, 由阶梯结构，$[Y_i,Y_j]$ 沿 $\partial_{u^1},\dots,\partial_{u^k}$ 的分量\emph{恒为零}：$\partial_{u^i}$ 与 $\partial_{u^j}$ 对易，剩余交叉项只产生沿 $\partial_{u^a}\,(a>k)$ 的分量。但 $\Delta$ 对合，故 $[Y_i,Y_j]\in\Delta=\operatorname{span}\{Y_1,\dots,Y_k\}$；而 $\Delta$ 中元素 $\sum_l\lambda^l Y_l$ 的第 $l\,(\le k)$ 个分量恰为 $\lambda^l$。比较两边前 $k$ 个分量得 $\lambda^l=0\,(\forall l)$，于是
\[
  [Y_i,Y_j]=\mathbf 0,\qquad\forall\,i,j .
\]

\emph{第 3 步（同时拉直）.}\, 对易向量场的流可交换，令
\[
  \Phi(t^1,\dots,t^k)=\phi^{Y_1}_{t^1}\circ\cdots\circ\phi^{Y_k}_{t^k}(\mathbf p_0),
\]
配上横截方向的 $n-k$ 个参数构成局部坐标，则 $Y_i=\partial_{t^i}$，故 $\Delta=\operatorname{span}\{\partial_{t^1},\dots,\partial_{t^k}\}$，水平叶 $\{t^{k+1}=c,\dots,t^n=c\}$ 即积分流形。$\Delta$ 可积。\qed
\end{proof}

\begin{lemma}[余维一情形：一行可算的微分形式判据 $\omega\wedge\mathrm d\omega=0$]\label{lem:cd-frobenius-1form}
设单个处处非零的 $1$-形式 $\omega$ 定义余维一分布 $\Delta=\ker\omega$（$\dim\Delta=n-1$）。则
\[
  \Delta\ \text{对合}\iff \omega([X,Y])=0\ \ \forall\,X,Y\in\Gamma(\Delta)\iff \omega\wedge\mathrm d\omega=0 .
\]
\end{lemma}
\begin{proof}
$\Delta=\ker\omega$ 对合即 $[X,Y]\in\ker\omega$，亦即 $\omega([X,Y])=0$，第一个等价显然。由外导数的内蕴公式
\[
  \mathrm d\omega(X,Y)=X\big(\omega(Y)\big)-Y\big(\omega(X)\big)-\omega([X,Y]),
\]
对 $X,Y\in\ker\omega$ 有 $\omega(X)=\omega(Y)=0$，故 $\mathrm d\omega(X,Y)=-\omega([X,Y])$。于是“$\omega([X,Y])=0\ \forall X,Y\in\ker\omega$”等价于“$\mathrm d\omega|_{\ker\omega}=0$”。再证后者 $\iff\omega\wedge\mathrm d\omega=0$：取余标架使 $\omega=e^1$、$\ker\omega=\operatorname{span}\{e_2,\dots,e_n\}$，写 $\mathrm d\omega=\sum_{i<j}a_{ij}\,e^i\wedge e^j$。则 $\mathrm d\omega|_{\ker\omega}=0\iff a_{ij}=0\ (\forall\,2\le i<j)$，即 $\mathrm d\omega=e^1\wedge\big(\textstyle\sum_j a_{1j}e^j\big)=\omega\wedge\eta$；而 $\omega\wedge\mathrm d\omega=e^1\wedge\sum_{i<j}a_{ij}e^i\wedge e^j=\sum_{2\le i<j}a_{ij}\,e^1\wedge e^i\wedge e^j$（凡含 $i=1$ 的项因 $e^1\wedge e^1=0$ 消失），故同样 $\iff a_{ij}=0\ (\forall\,2\le i<j)$。三者等价。\qed
\end{proof}

\begin{insight}[判据落地：独轮车的 $\omega\wedge\mathrm d\omega\neq0$ 即非完整]\label{ins:cd-frobenius-unicycle}
\cref{lem:cd-frobenius-1form} 把 \cref{thm:cd-frobenius} 的多约束行列式条件 $\mathrm d\omega_i\wedge\omega_1\wedge\cdots\wedge\omega_m=0$ 在\emph{单约束}时化为一行可算的检验 $\omega\wedge\mathrm d\omega=0$。\cref{ex:cd-three} 的独轮车正是它：$\omega=\sin\theta\,\mathrm dx-\cos\theta\,\mathrm dy$ 给 $\mathrm d\omega=\mathrm d\theta\wedge(\cos\theta\,\mathrm dx+\sin\theta\,\mathrm dy)$，而 $\omega\wedge\mathrm d\omega=-\,\mathrm dx\wedge\mathrm dy\wedge\mathrm d\theta\neq0$。故 $\Delta=\ker\omega$ \emph{非}对合、约束\emph{不可积}——这与 \cref{ch:nonholonomic} 用李括号 $[\mathbf g_1,\mathbf g_2]$ 逸出速度分布所得结论完全一致，只是这里从\emph{对偶}（余分布、$1$-形式）一侧看：$-\omega([\mathbf g_1,\mathbf g_2])=\mathrm d\omega(\mathbf g_1,\mathbf g_2)\neq0$ 恰度量了括号逸出 $\Delta$ 的“量”。可积性（Frobenius）与可控性（Chow）的“一体两面”（\cref{ins:nh-twoface}）在这条等式里看得最清楚。
\end{insight}

\subsection{四元分类体系}
约束沿四条独立维度分类（\cref{tab:cd-classify}）。

\begin{table}[H]\centering\small
\caption{约束的四维分类}
\label{tab:cd-classify}
\begin{tabularx}{\linewidth}{p{0.20\linewidth}LLL}
\toprule
维度 & 类型 A & 类型 B & 区分依据 \\
\midrule
时间依赖 & scleronomic（不含 $t$） & rheonomic（显含 $t$） & 约束是否显含时间 \\
等式/不等式 & bilateral（$\boldsymbol\varphi=\mathbf 0$） & unilateral（$\boldsymbol\varphi\ge\mathbf 0$） & Signorini 互补 $\varphi\lambda=0$ \\
可积性 & holonomic & nonholonomic & Frobenius（\cref{thm:cd-frobenius}） \\
理想性 & 理想（力 $\perp$ 虚位移） & 非理想（含摩擦） & 约束力是否做虚功 \\
\bottomrule
\end{tabularx}
\end{table}

\begin{example}[三类经典约束]\label{ex:cd-three}
\textbf{球铰} $\mathbf r_A-\mathbf r_B=\mathbf 0$：3D 完整、双边、定常，约束力为方向任意的 3D 反力。\;
\textbf{足端非穿透} $z_{\text{foot}}\ge0$：单边完整，足在空中时约束不激活（$\lambda=0$），触地时 $z=0,\ \lambda\ge0$，构成\emph{Signorini 互补}
\begin{equation}
  z_{\text{foot}}\ge0,\quad \lambda\ge0,\quad z_{\text{foot}}\,\lambda=0.
  \label{eq:cd-signorini-foot}
\end{equation}
\textbf{独轮车} $\dot x\sin\theta-\dot y\cos\theta=0$：非完整。写 1-形式 $\omega=\sin\theta\,\mathrm dx-\cos\theta\,\mathrm dy$，则 $\mathrm d\omega=\cos\theta\,\mathrm d\theta\wedge\mathrm dx+\sin\theta\,\mathrm d\theta\wedge\mathrm dy$，化简得 $\mathrm d\omega\wedge\omega\neq0$，故\emph{不可积}。
\end{example}

\paragraph{Chow--Rashevsky：与 Frobenius 对偶的可达性。}
独轮车的两个速度方向 $\partial/\partial v,\ \partial/\partial\omega$ 的 Lie 括号生成第三个方向，由\textbf{Chow--Rashevsky 定理}\cite{lynch2017modern}：若分布 bracket-generating（$\mathrm{Lie}(g_1,\dots,g_m)(\mathbf q)=T_{\mathbf q}\mathcal Q$），则任两点可由水平曲线连接。它与 Frobenius 恰成对偶——
\begin{center}\small
Frobenius $\Rightarrow$ 可积 $\Rightarrow$ holonomic $\Rightarrow$ 困于子流形\quad\textbf{vs}\quad Chow $\Rightarrow$ 不可积 $\Rightarrow$ nonholonomic $\Rightarrow$ 可控到全空间.
\end{center}

\begin{pitfall}[概念误区：「nonholonomic = 到不了某些位置」]
新手以为约束限制运动就限制了可达位置集。实则 Chow--Rashevsky 告诉我们：bracket-generating 的非完整系统能到达配置空间任意点——约束限制的是\emph{瞬时速度方向}，不是\emph{可达位置}。独轮车不能侧移，但通过转弯可到达任意 $(x,y,\theta)$。
\end{pitfall}

分类直接决定求解工具（\cref{tab:cd-tool}）：本章后续主攻完整双边（DAE/KKT，\cref{sec:cd-kkt}）与单边带摩擦（LCP/SOCP，\cref{sec:cd-contact}）；非完整系统的\emph{控制设计}（chained form 镇定等）属规划/移动机器人专题，本章只保留其分类与判据。

\begin{table}[H]\centering\small
\caption{约束类型到求解工具的映射}
\label{tab:cd-tool}
\begin{tabular}{p{0.26\linewidth}p{0.24\linewidth}p{0.30\linewidth}}
\toprule
约束类型 & 数学工具 & 代表方法 \\
\midrule
完整双边 & DAE / KKT & Baumgarte / GGL / proximal \\
完整单边 & LCP / NCP & Stewart--Trinkle / Moreau \\
非完整 & Pfaffian + 力层方程 & d'Alembert--Lagrange \\
单边 + 摩擦 & SOCP / CCP & Newton（凸化）/ SAP \\
\bottomrule
\end{tabular}
\end{table}

% ════════════════════════════════════════════════════════════════
\section{DAE 形式化与微分指数}
\label{sec:cd-dae}

\paragraph{动机：为何直接积分约束系统会发散。}
自由系统 $\mathbf M\ddot{\mathbf q}+\mathbf C\dot{\mathbf q}+\mathbf g=\boldsymbol\tau$ 是 ODE，任何标准积分器都能正确处理。加上约束 $\boldsymbol\varphi(\mathbf q)=\mathbf 0$ 后变成
\begin{equation}
  \mathbf M(\mathbf q)\ddot{\mathbf q}+\mathbf C(\mathbf q,\dot{\mathbf q})\dot{\mathbf q}+\mathbf g(\mathbf q)=\boldsymbol\tau+\mathbf J^\top(\mathbf q)\boldsymbol\lambda,\qquad \boldsymbol\varphi(\mathbf q)=\mathbf 0,
  \label{eq:cd-dae}
\end{equation}
它\emph{不再是 ODE}：既含代数约束 $\boldsymbol\varphi=\mathbf 0$，又含未知约束力 $\boldsymbol\lambda$——这是一个\textbf{微分代数方程（DAE）}。若忽略约束、只积分动力学部分而「指望」约束自然满足，后果是灾难性的。

\subsection{微分指数与约束力学的指数阶梯}
\begin{definition}[微分指数\cite{hairer1996solving}]\label{def:cd-index}
对半显式 DAE $\dot{\mathbf x}=\mathbf f(\mathbf x,\mathbf y,t),\ \mathbf 0=\mathbf g(\mathbf x,\mathbf y,t)$，其\emph{微分指数} $\nu$ 是把代数约束 $\mathbf g$ 关于 $t$ 微分到能\emph{显式}解出代数变量 $\mathbf y$ 所需的最小次数再加 $1$。
\end{definition}
\noindent 直觉：微分指数衡量「代数约束藏得有多深」。index-1 表示对 $\mathbf g=\mathbf 0$ 微分一次就看见 $\mathbf y$；index-3 表示要微分三次。每多一阶，数值求解就难一截。对机器人约束 $\boldsymbol\varphi(\mathbf q)=\mathbf 0$，按写成的层级不同，指数不同（\cref{tab:cd-ladder}）。

\begin{table}[H]\centering\small
\caption{约束力学的指数阶梯}
\label{tab:cd-ladder}
\begin{tabular}{p{0.16\linewidth}p{0.42\linewidth}p{0.10\linewidth}p{0.20\linewidth}}
\toprule
层级 & 方程 & 指数 & 说明 \\
\midrule
位置层 & $\mathbf M\ddot{\mathbf q}=\boldsymbol\tau+\mathbf J^\top\boldsymbol\lambda,\ \boldsymbol\varphi(\mathbf q)=\mathbf 0$ & $3$ & 原始形式，最病态 \\
速度层 & $\mathbf M\ddot{\mathbf q}=\boldsymbol\tau+\mathbf J^\top\boldsymbol\lambda,\ \mathbf J\dot{\mathbf q}=\mathbf 0$ & $2$ & 对 $\boldsymbol\varphi$ 微分一次 \\
加速度层 & $\mathbf M\ddot{\mathbf q}=\boldsymbol\tau+\mathbf J^\top\boldsymbol\lambda,\ \mathbf J\ddot{\mathbf q}+\dot{\mathbf J}\dot{\mathbf q}=\mathbf 0$ & $1$ & 对 $\boldsymbol\varphi$ 微分两次 \\
约化坐标 & $\ddot{\mathbf z}=\mathbf F(\mathbf z,\dot{\mathbf z})$ on $T_{\mathbf q}\mathcal M$ & $0$ & 完全消去约束 \\
\bottomrule
\end{tabular}
\end{table}

\begin{insight}[index-3 之源：差两阶微分]
约束力 $\boldsymbol\lambda$ \emph{本质上是加速度层的量}（它平衡力），而约束条件 $\boldsymbol\varphi=\mathbf 0$ 是位置层的量。两者之间差了\emph{两阶微分}——这个「阶差」正是 index-3 的来源。理解这一点，就理解了 drift-off 与一切稳定化方法的根本动机。
\end{insight}

\subsection{Drift-off 效应及其 \texorpdfstring{$O(t^2)$}{O(t2)} 增长}
\begin{theorem}[drift-off，Petzold 1982\cite{petzold1982dae,hairer1996solving}]\label{thm:cd-drift}
对 index-3 DAE，标准 BDF 方法的阶数降低 $\nu-1=2$ 阶；约束违反在速度层线性增长、位置层二次增长。
\end{theorem}
\begin{proof}[骨架证明]
设 $m=1,\ \mathbf M=\mathbf I$，$\mathbf J$ 为常向量，忽略非保守力。令 $e(t):=\boldsymbol\varphi(\mathbf q(t))$ 为约束违反，真实解 $e\equiv0$。在 index-1（加速度层）形式下，方程\emph{只}强制 $\ddot e=\mathbf J\ddot{\mathbf q}+\dot{\mathbf J}\dot{\mathbf q}=0$，而 $\dot e=\mathbf J\dot{\mathbf q}$ 与 $e=\boldsymbol\varphi$ 本身\emph{未被方程约束}。离散积分引入 $O(h^p)$ 局部截断误差 $\varepsilon$，于是
\begin{equation}
  \ddot e=\varepsilon(t)\ \Longrightarrow\ \dot e(t)\sim\varepsilon\,t,\quad e(t)\sim\varepsilon\,t^2.
  \label{eq:cd-driftoff}
\end{equation}
即速度层漂移线性增长、位置层漂移二次增长。\qed
\end{proof}

\begin{insight}[火车与铁轨]
把约束流形想成铁轨，数值解是行驶其上的火车。每一步都有微小侧向力（截断误差）。没有「拉回铁轨」的机制，火车就越偏越远——这就是 drift-off。Baumgarte 稳定化（\cref{sec:cd-baumgarte}）相当于给铁轨加了弹簧与阻尼器。
\end{insight}

\paragraph{指数约化的四条策略。}
\textbf{(1) 微分到 index-1}：把约束微分两次得加速度层，可用标准 ODE 积分器，但丢失位置/速度层信息、产生 drift-off。\;
\textbf{(2) 超定 DAE（OD-DAE）}：同时保留三层方程，用 DASSL 类专用求解器。\;
\textbf{(3) 约化坐标}：取满足约束的最小坐标 $\mathbf z\in\mathbb R^{n-m}$ 使 $\mathbf q=\boldsymbol\phi(\mathbf z)$ 自动满足约束，变回纯 ODE，但 $\boldsymbol\phi$ 可能难构造或有奇异。\;
\textbf{(4) 稳定化}：保持 index-1 并加稳定化项抑制 drift-off——工程最常用，即 \cref{sec:cd-baumgarte}。

\begin{pitfall}[思维陷阱：「用更高阶积分器就能消除 drift-off」]
新手以为 RK4 阶太低、换 RK8/DOP853 即可。但 \cref{thm:cd-drift} 表明：对 index-3 DAE，\emph{任何}标准 ODE 积分器都会 drift-off，只是速率不同。高阶方法只让初始偏离更小（$O(h^p)$ 而非 $O(h)$），\emph{增长趋势不变}。必须用约束稳定化或专用 DAE 求解器。
\end{pitfall}

% ════════════════════════════════════════════════════════════════
\section{Lagrange 乘子法与 KKT 系统}
\label{sec:cd-kkt}

\paragraph{动机：约束力从哪里来。}
\cref{ch:geometric-mechanics} 用虚功原理推无约束系统时，虚位移 $\delta\mathbf q$ 可取任意方向，故能从 $\sum(\mathbf F_i-m_i\ddot{\mathbf r}_i)\cdot\delta\mathbf r_i=0$ 推出每个分量独立为零。有约束时\emph{完全不同}：虚位移必须满足约束线性化 $\mathbf J\delta\mathbf q=\mathbf 0$，于是只能推出「广义力在约束切空间方向为零」，而约束\emph{法向}的力是自由的、由约束面自动提供——这就是约束力。

\subsection{从 d'Alembert 到 KKT 的每一步}
\textbf{第 1 步（d'Alembert--Lagrange 原理）.} 用广义坐标重写虚功原理，对满足约束的虚位移有
\begin{equation}
  \Big(\tfrac{\mathrm d}{\mathrm dt}\tfrac{\partial L}{\partial\dot{\mathbf q}}-\tfrac{\partial L}{\partial\mathbf q}-\boldsymbol\tau\Big)\cdot\delta\mathbf q=0,\qquad\forall\,\delta\mathbf q\in\ker\mathbf J(\mathbf q).
  \label{eq:cd-dalembert}
\end{equation}
\textbf{第 2 步（Lagrange 乘子定理）.} 关键线性代数：\emph{若 $\mathbf Q^\top\boldsymbol\xi=0$ 对一切 $\boldsymbol\xi\in\ker\mathbf J$ 成立，则 $\exists\,\boldsymbol\lambda\in\mathbb R^m$ 使 $\mathbf Q=\mathbf J^\top\boldsymbol\lambda$}。这正是 $(\ker\mathbf J)^\perp=\mathrm{row}(\mathbf J)=\mathrm{col}(\mathbf J^\top)$ 的推论（\cref{ch:matderiv}）。物理含义：对一切容许虚位移做功为零的力，必属于约束 Jacobian 的行空间——即它是约束力。\;
\textbf{第 3 步（EL + 约束力）.} 由此 $\tfrac{\mathrm d}{\mathrm dt}\tfrac{\partial L}{\partial\dot{\mathbf q}}-\tfrac{\partial L}{\partial\mathbf q}=\boldsymbol\tau+\mathbf J^\top\boldsymbol\lambda$，代入机器人标准形得 \cref{eq:cd-dae} 的力层方程。\;
\textbf{第 4 步（约束两次微分）.} $\boldsymbol\varphi=\mathbf 0\Rightarrow\mathbf J\dot{\mathbf q}=\mathbf 0\Rightarrow\mathbf J\ddot{\mathbf q}+\dot{\mathbf J}\dot{\mathbf q}=\mathbf 0$。\;
\textbf{第 5 步（联立得 KKT 系统）.}
\begin{equation}
  \boxed{\;\begin{bmatrix}\mathbf M(\mathbf q) & -\mathbf J^\top(\mathbf q)\\ \mathbf J(\mathbf q) & \mathbf 0\end{bmatrix}\begin{bmatrix}\ddot{\mathbf q}\\ \boldsymbol\lambda\end{bmatrix}=\begin{bmatrix}\boldsymbol\tau-\mathbf C\dot{\mathbf q}-\mathbf g\\ -\dot{\mathbf J}\dot{\mathbf q}\end{bmatrix}\;}
  \label{eq:cd-kkt}
\end{equation}

\subsection{KKT 系统的结构分析}
\cref{eq:cd-kkt} 是一个\textbf{鞍点问题}：左上块 $\mathbf M\succ\mathbf 0$ 正定，$\pm\mathbf J^\top$ 耦合动力学与约束，整体矩阵\emph{不定}（既有正负特征值），但在 LICQ 下唯一可解。\textbf{LICQ}（线性无关约束规范）即 $\mathbf J$ 行满秩；违反时（冗余约束或奇异构型）$\mathbf G=\mathbf J\mathbf M^{-1}\mathbf J^\top$ 奇异，需正则化。

\begin{theorem}[Schur 消元与 Delassus 矩阵\cite{khatib1987unified,featherstone2008rigid}]\label{thm:cd-schur}
设 $\mathbf G:=\mathbf J\mathbf M^{-1}\mathbf J^\top$（\emph{Delassus 矩阵}，约束空间逆惯量），其逆 $\boldsymbol\Lambda:=\mathbf G^{-1}$ 即 Khatib(1987) \emph{操作空间惯量矩阵}。由 \cref{eq:cd-kkt} 第一行解出 $\ddot{\mathbf q}=\mathbf M^{-1}(\boldsymbol\tau-\mathbf C\dot{\mathbf q}-\mathbf g+\mathbf J^\top\boldsymbol\lambda)$，代入第二行得闭式
\begin{equation}
  \boldsymbol\lambda=-\mathbf G^{-1}\big[\mathbf J\mathbf M^{-1}(\boldsymbol\tau-\mathbf C\dot{\mathbf q}-\mathbf g)+\dot{\mathbf J}\dot{\mathbf q}\big],\qquad \ddot{\mathbf q}=\mathbf M^{-1}(\boldsymbol\tau-\mathbf C\dot{\mathbf q}-\mathbf g+\mathbf J^\top\boldsymbol\lambda).
  \label{eq:cd-schur}
\end{equation}
\end{theorem}

\begin{note}
\textbf{记号订正（与 N/O 章统一）.}\;
当代约束/接触动力学文献（Carpentier、Featherstone--Orin Handbook、Pinocchio）中，\emph{Delassus 矩阵/算子标准指 $\mathbf G=\mathbf J\mathbf M^{-1}\mathbf J^\top$ 本身}（名自 É.~Delassus 1920 年代的滑动摩擦/约束力工作），其\emph{逆} $\boldsymbol\Lambda=\mathbf G^{-1}$ 才是 Khatib 1987 的操作空间惯量。本书全章遵此：$\mathbf G$ 是 Delassus、$\boldsymbol\Lambda$ 是 OSIM，二者互逆而非等同。
\end{note}

\paragraph{$\mathbf J$ 的几何与 $\boldsymbol\lambda$ 的物理。}
约束 Jacobian $\mathbf J=\partial\boldsymbol\varphi/\partial\mathbf q$ 的\emph{行空间} $\mathrm{col}(\mathbf J^\top)$ 是约束力方向（$\mathbf J^\top\boldsymbol\lambda$ 只能产生此方向的力），\emph{零空间} $\ker\mathbf J$ 是容许运动空间即约束流形切空间 $T_{\mathbf q}\mathcal M$；二者在 M-内积下正交互补 $\mathbb R^{n}=\ker(\mathbf J)\oplus_{\mathbf M}\mathrm{range}(\mathbf M^{-1}\mathbf J^\top)$。乘子 $\boldsymbol\lambda$ 的物理维度随约束而定（\cref{tab:cd-lambda}）。当 $\boldsymbol\tau=\mathbf C\dot{\mathbf q}+\mathbf g$ 且 $\dot{\mathbf J}\dot{\mathbf q}=\mathbf 0$ 时 $\boldsymbol\lambda=\mathbf 0$——不需额外约束力；有外力试图违反约束时，$\boldsymbol\lambda$ 恰好产生补偿力把系统拉回流形，不多不少（这正是 Gauss 原理的精神，\cref{sec:cd-gauss}）。

\begin{table}[H]\centering\small
\caption{乘子 $\boldsymbol\lambda$ 的物理含义与维度}
\label{tab:cd-lambda}
\begin{tabular}{p{0.30\linewidth}p{0.42\linewidth}p{0.12\linewidth}}
\toprule
约束类型 & $\boldsymbol\lambda$ 物理含义 & 维度 \\
\midrule
球铰（点对点） & 3D 反力 & $3$ \\
转动铰（revolute） & 5D wrench（除旋转轴外） & $5$ \\
足端点接触 & 3D 接触力 $(f_x,f_y,f_z)$ & $3$ \\
足端面接触 & 6D wrench（力 + 力矩） & $6$ \\
闭链内力 & cut joint 处 wrench & $5$ 或 $6$ \\
\bottomrule
\end{tabular}
\end{table}

\begin{pitfall}[编程陷阱与 LICQ 误解]
\textbf{(转置)} 写成 $\mathbf M\ddot{\mathbf q}+\mathbf h=\boldsymbol\tau+\mathbf J\boldsymbol\lambda$（缺转置）会让力落在错误空间——约束力\emph{永远}是 $\mathbf J^\top\boldsymbol\lambda$。\;
\textbf{(LICQ)} 「$\mathbf J$ 不满秩 = 模型有 bug」是误解：四足 $4\times3=12$ 约束、只需 $6$ 维控质心，冗余是常态；正确处理是 proximal 正则化 $(\mathbf G+\mu\mathbf I)^{-1}$ 或伪逆（\cref{sec:cd-gauss}）。
\end{pitfall}

% ════════════════════════════════════════════════════════════════
\section{Baumgarte 稳定化与替代方案}
\label{sec:cd-baumgarte}

\paragraph{动机：index-1 形式不够。}
\cref{sec:cd-dae} 表明加速度层（index-1）只强制 $\ddot e=0$，故即便 $e(0)\neq0$ 或 $\dot e(0)\neq0$ 也不被修正——误差被「冻结」或线性增长。工程需要一个\emph{简单、廉价、可调}的方法抑制漂移。Baumgarte(1972) 在 \emph{Computer Methods in Applied Mechanics and Engineering} 创刊卷上提出一个优雅思路：\textbf{把约束违反当作被微扰的动力系统，加负反馈使其稳定}。半个世纪后它仍是 MuJoCo/Pinocchio/Drake 的核心组件之一。

\subsection{完整推导：误差方程与临界阻尼}
定义约束违反 $e(t):=\boldsymbol\varphi(\mathbf q(t))$。把加速度层方程 $\mathbf J\ddot{\mathbf q}+\dot{\mathbf J}\dot{\mathbf q}=\mathbf 0$ 替换为加了阻尼项 $2\alpha\mathbf J\dot{\mathbf q}$ 与弹性项 $\beta^2\boldsymbol\varphi$ 的形式：
\begin{equation}
  \mathbf J\ddot{\mathbf q}+\dot{\mathbf J}\dot{\mathbf q}+2\alpha\,\mathbf J\dot{\mathbf q}+\beta^2\,\boldsymbol\varphi(\mathbf q)=\mathbf 0.
  \label{eq:cd-baumgarte-feedback}
\end{equation}
由 $e=\boldsymbol\varphi,\ \dot e=\mathbf J\dot{\mathbf q},\ \ddot e=\mathbf J\ddot{\mathbf q}+\dot{\mathbf J}\dot{\mathbf q}$ 代入，得二阶线性误差 ODE 及其特征方程
\begin{equation}
  \ddot e+2\alpha\,\dot e+\beta^2 e=0,\qquad s^2+2\alpha s+\beta^2=0,\qquad s_{1,2}=-\alpha\pm\sqrt{\alpha^2-\beta^2}.
  \label{eq:cd-baumgarte-error}
\end{equation}
这正是一个弹簧-阻尼系统。阻尼分类：$\alpha<\beta$ 欠阻尼（复根、振荡衰减）；$\alpha>\beta$ 过阻尼（双实根、最慢模越大越慢）；\textbf{$\alpha=\beta$ 临界阻尼}（双重根 $s=-\alpha$，无振荡最快衰减 $e(t)=(c_1+c_2t)e^{-\alpha t}$）。故工程取 $\alpha=\beta$。

\begin{insight}[Baumgarte 把约束违反当弹簧-阻尼]
\cref{eq:cd-baumgarte-feedback} 的两个新增项无任何物理来源，纯是\emph{反馈控制律}：弹性项 $\beta^2\boldsymbol\varphi$ 像弹簧把违反拉回零，阻尼项 $2\alpha\mathbf J\dot{\mathbf q}$ 像阻尼器抑制振荡。把约束稳定化看成一个二阶伺服回路，临界阻尼、参数整定、稳定域全都是控制论的常识。
\end{insight}

\subsection{参数准则与缺陷}
\textbf{准则 1（时间常数）}：欲让违反在 $T_{\text{settle}}$ 内衰到 $1\%$，由 $e^{-\alpha T_{\text{settle}}}\approx0.01$ 得 $\alpha\approx4.6/T_{\text{settle}}$。\;
\textbf{准则 2（步长约束）}：显式 Euler 稳定域要求 $\alpha h\le1$（Flores, EUCOMES 2008）。\;
\cref{tab:cd-baumgarte} 给经验范围。Baumgarte 有四个固有缺陷：\textbf{(1) 破坏辛结构}——反馈项不来自任何 Hamilton 量，添加后不再保能量（\cref{ch:geometric-mechanics} 辛积分器的能量有界优势被破坏）；\textbf{(2) 参数依赖步长}——不存在对所有步长都最优的 $(\alpha,\beta)$，换步长就得换参数；\textbf{(3) $\mathbf J$ 奇异时失效}——$\beta^2\boldsymbol\varphi$ 的修正方向（即 $\mathbf J$）退化；\textbf{(4) 漂移仅被抑制非消除}——稳态误差 $e_{\text{ss}}\sim F/(\alpha^2\beta^2)$，永不精确回到约束面。

\begin{table}[H]\centering\small
\caption{Baumgarte 参数经验范围（取 $\alpha=\beta$ 临界阻尼）}
\label{tab:cd-baumgarte}
\begin{tabular}{p{0.30\linewidth}p{0.16\linewidth}p{0.22\linewidth}p{0.22\linewidth}}
\toprule
应用场景 & 步长 $h$ & 推荐 $\alpha=\beta$ & 理由 \\
\midrule
腿足 MPC（1\,kHz） & 1\,ms & $100$–$500$\,Hz & 快修正、不超稳定域 \\
仿真（2\,kHz，隐式） & 0.5\,ms & $200$–$1000$\,Hz & 隐式可承受更高 \\
慢速工业臂（100\,Hz） & 10\,ms & $10$–$50$\,Hz & 步长大需保守 \\
\bottomrule
\end{tabular}
\end{table}

\begin{pitfall}[思维陷阱：「Baumgarte 越强越好」]
新手以为 $\alpha=10000$ 收敛最快。但当 $\alpha h>1$ 时显式积分器\emph{发散}——$\alpha$ 太大反而让约束违反爆炸。取 $\alpha\le1/(h\cdot\text{安全系数})$，安全系数 $2$–$3$。
\end{pitfall}

\subsection{四条替代方案}
Baumgarte 的缺陷催生了更稳健的替代：

\begin{itemize}
  \item \textbf{坐标投影（Eich 1993）}：每步积分后解 $\min_{\tilde{\mathbf q}}\|\tilde{\mathbf q}-\mathbf q\|_{\mathbf M}^2\ \text{s.t.}\ \boldsymbol\varphi(\tilde{\mathbf q})=\mathbf 0$，用 Newton 把 $\mathbf q$ 拉回流形——\emph{严格消除漂移}，代价是每步 $1$–$2$ 次 Newton 迭代。
  \item \textbf{GGL（Gear--Gupta--Leimkuhler 1985）}\cite{gear1985automatic}：引入第二乘子 $\boldsymbol\eta$，同时强制位置与速度约束 $\dot{\mathbf q}=\mathbf v+\mathbf J^\top\boldsymbol\eta,\ \mathbf M\dot{\mathbf v}=\mathbf f+\mathbf J^\top\boldsymbol\lambda,\ \boldsymbol\varphi=\mathbf 0,\ \mathbf J\dot{\mathbf q}=\mathbf 0$，把 drift-off 阶降为零，代价是系统维度增加（稳定化 index-2 形式）。
  \item \textbf{SHAKE(1977)/RATTLE(1983)}\cite{ryckaert1977numerical,andersen1983rattle}：\emph{辛} DAE 积分器。SHAKE 在 Verlet 步后解非线性方程把位置拉回约束面；RATTLE 是其速度层扩展，同时满足位置与速度约束。\emph{辛 $+$ 约束满足}使其成为分子动力学（GROMACS/LAMMPS）标准——为何优于 Baumgarte：它\emph{保}辛而 Baumgarte \emph{破}辛（cref \cref{ch:geometric-mechanics} 辛积分器）。
  \item \textbf{后稳定化（Ascher--Chin--Petzold--Reich 1995）}：与积分器独立的正交投影稳定器，先用任意积分器前进一步再投影，可叠加使用、不改积分器本身。
\end{itemize}

\begin{insight}[MuJoCo 软约束 = 软化的 Baumgarte]
MuJoCo（Todorov 2014）不直接加约束力，而解凸优化 $\min_{\mathbf x}\tfrac12\|\mathbf x-\mathbf a_{\text{free}}\|_{\mathbf M}^2+\tfrac12\|\mathbf y-\mathbf a^*\|_{\mathbf R^{-1}}^2$，其中 $\mathbf a^*=-2\alpha\mathbf v-\beta^2\mathbf r$、$\mathbf R=(1-d)/d\cdot\mathbf G$。当 $\mathbf R\to\mathbf 0$（硬约束）即退回临界阻尼 Baumgarte。这把「Baumgarte 反馈」和「软接触正则化」统一成同一对象的两端，是 \cref{sec:cd-diffcontact} 可微接触的伏笔。\textbf{注意} MuJoCo 的 \texttt{solref}$=(1/\beta,\ \alpha/\beta)$ 是 (timeconst, dampratio) 而非 $(\alpha,\beta)$，易混。
\end{insight}

% ════════════════════════════════════════════════════════════════
\section{闭链机构}
\label{sec:cd-loop}

\paragraph{动机：树形模型的局限。}
URDF 与 \cref{ch:spatial-dynamics} 的 $O(N)$ 算法都假设机器人是\emph{树形}——从根到每个末端只有一条路径。但许多实际机构有\emph{闭合运动链}：Stewart--Gough 平台（6-UPS 并联）、Delta 机器人（$3$ 个平行四边形 $+$ 平台 loop）、Cassie/Digit 每腿一个四杆联动（胫骨-跟腱连杆-足跟弹簧）、人形四杆膝。它们的 URDF 表示必须「砍掉一个关节变成树 $+$ 加回环路闭合约束」，否则 ABA 会\emph{忽略闭环耦合力}、算出错误动力学。

\begin{insight}[闭链像桁架]
闭链约束像桥梁桁架：只算一侧（树形）、忽略另一侧支撑力，桥就「塌」了。cut-joint 方法相当于「切开桁架但记住切口处有内力」——内力就是闭链约束力 $\boldsymbol\lambda$。
\end{insight}

\subsection{建模流程}
沿 Featherstone 2008 第 8 章\cite{featherstone2008rigid}：

\begin{enumerate}
  \item \textbf{识别 loop}：由 Grübler 公式 $\mathrm{DOF}=6(N-1-j)+\sum_{i=1}^j f_i$（$N$ 含地面的刚体数，$j$ 关节数，$f_i$ 关节自由度）算出环路数。
  \item \textbf{选 cut joint}：每个 loop 切一个关节使剩余成 spanning tree；原则是优先切高自由度关节（减少约束维度）、避免切驱动关节。
  \item \textbf{环路闭合约束}：被切关节连接刚体 $A,B$，要求经 spanning tree 的两条路径到达切口两端的 $\SE(3)$ 变换相等
  \begin{equation}
    \mathbf C_i(\mathbf q)={}^0\mathbf M_{A_i}^{-1}\,{}^0\mathbf M_{B_i}=\mathbf I_{4\times4}.
    \label{eq:cd-loopclosure}
  \end{equation}
  \item \textbf{线性化得 Pfaffian}：$\mathbf K_i(\mathbf q)\dot{\mathbf q}=\mathbf 0$，$\mathbf K_i$ 为环路闭合约束 Jacobian。
  \item \textbf{tree 正动力学 + 约束力}：设 ABA 算出的「无约束加速度」为 $\hat{\mathbf a}$，由 Schur 补解
  \begin{equation}
    \mathbf K\mathbf M^{-1}\mathbf K^\top\boldsymbol\lambda=-\mathbf K\hat{\mathbf a}-\dot{\mathbf K}\dot{\mathbf q},\qquad \ddot{\mathbf q}=\hat{\mathbf a}+\mathbf M^{-1}\mathbf K^\top\boldsymbol\lambda.
    \label{eq:cd-loop-schur}
  \end{equation}
\end{enumerate}

切口约束维度按关节类型定：3D 球铰 $3$、转动铰 $5$（$6$ 减 $1$ 旋转）、固连 $6$、棱柱 $5$。以 Cassie 为例，每腿四杆联动取 spanning tree 髋 $\to$ 大腿 $\to$ 胫骨 $\to$ 足、cut joint 取跟腱连杆与足的连接点（5D 转动铰闭合）；整机 $n=20$ 树关节 $+\ 2\times5=10$ 维闭链约束，KKT 矩阵 $30\times30$。

\subsection{O(N) 求解复杂度}
\cref{eq:cd-loop-schur} 的关键是\emph{先 $O(N)$ 的 ABA 解树部分、再 Schur 补解闭链约束力}，总复杂度
\begin{equation}
  O(N)+O(m^3),
  \label{eq:cd-loop-complexity}
\end{equation}
其中 $N$ 树自由度、$m$ 约束维度。由于通常 $m\ll N$，这远快于稠密 KKT 的 $O((N+m)^3)$。Sathya--Carpentier 的 \textbf{constrainedABA}\cite{sathya2024constrained}（IEEE TRO 2024）实现了这一线性复杂度策略，并对冗余/奇异约束保持稳健；对 Atlas 这类约 $36$ 自由度机器人，约束前向动力学可达\emph{微秒级}，比稠密 KKT 快数倍。

\begin{pitfall}[编程陷阱：URDF 表达不了闭链]
试图在 URDF 用 \texttt{mimic} 关节模拟闭链是错的——\texttt{mimic} 只是运动学耦合，\emph{不产生约束力}。正确做法是用支持环路闭合的格式（如 MuJoCo XML）或手工构造刚性约束模型（Pinocchio \texttt{RigidConstraintModel}）。\cref{ch:control} 机械臂部并联臂将特化本节框架到 Delta/Stewart 闭链。
\end{pitfall}

% ════════════════════════════════════════════════════════════════
\section{接触动力学}
\label{sec:cd-contact}

\paragraph{动机：接触为何比一般约束难得多。}
闭链约束是\emph{双边}（始终激活）的。足端接触却是\emph{单边}（可能激活也可能不激活）$+$\emph{带摩擦}的，带来三个根本困难：\textbf{(1) 互补性} $z\ge0,\ \lambda\ge0,\ z\lambda=0$——不是连续优化而是组合问题；\textbf{(2) 摩擦锥} $\|\mathbf f_t\|\le\mu_f f_n$——非线性不等式；\textbf{(3) 事件切换}——接触建立/释放/滑动/粘着的切换是混合系统。

\begin{insight}[接触本质上不是连续问题]
接触动力学之难，不在方程复杂，而在它\emph{本质上不是连续问题}：解的存在唯一性在一般情况下都不保证（Painlevé 悖论 1895\cite{painleve1895lois}：刚体 $+$ Coulomb 摩擦可无解/多解/不一致）。一切「求解接触」的方法都是某种\emph{近似}，区别只在近似方式与代价。
\end{insight}

\begin{note}
\textbf{与 N 章的分工.}\;
本节把接触放在\emph{约束动力学统一框架}内——它是 unilateral 约束特例，与 DAE/KKT/Baumgarte/IFT 的连接是本章主旨。摩擦锥/GIWC 的几何、LCP/CCP 求解器、接触力分配的算法与工程主体在 \cref{ch:contact-mechanics}（接触力学专论）。二者互引，Signorini/Stewart--Trinkle/Moreau 不重复推导。
\end{note}

\subsection{Signorini 条件与 Coulomb 摩擦锥}
设 $\phi_n$ 为法向间距（gap function）、$\lambda_n$ 为法向接触力，\textbf{Signorini 条件}
\begin{equation}
  \phi_n\ge0,\qquad \lambda_n\ge0,\qquad \phi_n\,\lambda_n=0
  \label{eq:cd-signorini}
\end{equation}
三条分别表示：不穿透（$\phi_n\ge0$）、法向力只能是压力（$\lambda_n\ge0$，地面不能「拉住」脚）、互补（$\phi_n\lambda_n=0$：要么不接触有间距无力、要么接触无间距有力，不能既穿透又有力）。\textbf{Coulomb 摩擦锥}要求接触力 $\mathbf f=(\mathbf f_t,f_n)$ 满足
\begin{equation}
  \|\mathbf f_t\|\le\mu_f\,f_n,
  \label{eq:cd-coulomb}
\end{equation}
几何上是一个「冰激凌锥」。工程常用 $k$-边形线性化：$|f_{t,x}|+|f_{t,y}|\le\mu_f f_n$（4-pyramid），或选 $k$ 个方向 $\mathbf d_i$ 令 $\mathbf d_i^\top\mathbf f_t\le\mu_f f_n$（$k$-polygon）。

\begin{pitfall}[概念误区：「4 边摩擦锥就够」]
4-pyramid 在对角方向（$45^\circ$）允许的摩擦力比圆锥少 $1-1/\sqrt2\approx29\%$。对大摩擦系数（橡胶足 $\mu_f=0.8$）影响显著，建议至少 8-polygon 或直接用 SOCP。
\end{pitfall}

\subsection{LCP/NCP/CCP 表述}
在加速度层，法向加速度 $a_n=\mathbf J_n\ddot{\mathbf q}+\dot{\mathbf J}_n\dot{\mathbf q}$，代入动力学得 $a_n=\mathbf J_n\mathbf M^{-1}\mathbf J_n^\top\boldsymbol\lambda_n+\mathbf J_n\mathbf M^{-1}(\boldsymbol\tau-\mathbf h)+\dot{\mathbf J}_n\dot{\mathbf q}$。令 $\mathbf A=\mathbf J_n\mathbf M^{-1}\mathbf J_n^\top$（\emph{又是 Delassus 矩阵 $\mathbf G$}）、$\mathbf b=\mathbf J_n\mathbf M^{-1}(\boldsymbol\tau-\mathbf h)+\dot{\mathbf J}_n\dot{\mathbf q}$，Signorini 化为\textbf{线性互补问题（LCP）}
\begin{equation}
  0\le\boldsymbol\lambda_n\ \perp\ (\mathbf A\boldsymbol\lambda_n+\mathbf b)\ge0.
  \label{eq:cd-lcp}
\end{equation}
加入 Coulomb 摩擦后切向方程非线性（最大耗散原则：滑动时 $\mathbf f_t=-\mu_f f_n\,\mathbf v_t/\|\mathbf v_t\|$、粘着时 $\|\mathbf f_t\|\le\mu_f f_n$），给出\textbf{非线性互补问题（NCP）}或\textbf{锥互补问题（CCP）}。

\subsection{时步法：Stewart--Trinkle 与 Moreau}
加速度层 LCP 有两个根本缺陷：Painlevé 悖论（某些构型无解/多解）、无法处理速度不连续（impact）。\textbf{时步法}绕过这两点——不在加速度层而在\emph{冲量层}（离散时间步）求解。

\begin{theorem}[Stewart--Trinkle 时步法 1996\cite{stewart1996implicit}]\label{thm:cd-stewart}
设步长 $h$，离散动力学
\begin{equation}
  \mathbf M(\mathbf v^+-\mathbf v^-)=h\,\mathbf f_{\text{ext}}+\mathbf J_n^\top\mathbf p_n+\mathbf J_t^\top\mathbf p_t,
  \label{eq:cd-stewart}
\end{equation}
其中 $\mathbf p_n,\mathbf p_t$ 是法向/切向接触\emph{冲量}（$\mathbf p=h\boldsymbol\lambda$）。配合离散 Signorini $0\le\mathbf p_n\perp\boldsymbol\phi_n^+\ge0$ 与离散 Coulomb $\mathbf p_t\in-\mu_f\mathbf p_n\,\partial\|\mathbf v_t^+\|$，整体形成 mixed LCP / CCP，在 $h\to0$ 时其离散解收敛到\emph{measure differential inclusion（MDI）}的解。
\end{theorem}
\noindent 优势：避免加速度层的 Painlevé 病态（冲量层无此问题）、\emph{统一}持续接触与碰撞、存在唯一性保证更强。

\textbf{Moreau 半隐扫过过程 1988}\cite{moreau1988unilateral} 把接触视为 MDI $\mathbf M\,\mathrm d\mathbf v\in-\mathbf h\,\mathrm dt+N_K(\mathbf v)\,\mathrm dt$（$N_K$ 为容许速度集 $K$ 的法锥），三步推进：预测 $\tilde{\mathbf v}=\mathbf v_k+h\mathbf M^{-1}\mathbf f_{\text{ext}}$、修正 $\mathbf v_{k+1}=\mathrm{proj}_{K(\mathbf q_k)}^{\mathbf M}(\tilde{\mathbf v})$（M-范数投影到容许集）、更新 $\mathbf q_{k+1}=\mathbf q_k+h\mathbf v_{k+1}$。它无需事件检测（自动处理 make/break/slide/stick）、一步一解、凸投影总有解。

\textbf{Anitescu 凸化 2006}\cite{anitescu2006optimization} 把 Coulomb 时步问题凸松弛为\emph{严格凸} QP/CCP $\min_{\mathbf v^+}\tfrac12(\mathbf v^+-\tilde{\mathbf v})^\top\mathbf M(\mathbf v^+-\tilde{\mathbf v})\ \text{s.t.}\ \Phi_i(\mathbf v^+)\in\mathcal{FC}_i^*$：滑动时精确、粘着时有随 $h\to0$ 消失的微小人工滑移。它是 MuJoCo/Dojo/Drake/Chrono 的\emph{共同凸根基}。\cref{tab:cd-contact-events} 给出接触事件的完整分类。

\begin{table}[H]\centering\small
\caption{接触事件分类}
\label{tab:cd-contact-events}
\begin{tabular}{p{0.20\linewidth}p{0.22\linewidth}p{0.26\linewidth}p{0.20\linewidth}}
\toprule
事件 & 物理 & 数学条件 & 处理 \\
\midrule
make（建立） & 物体到达地面 & $\phi\to0^+,\ \dot\phi<0$ & 激活约束行 \\
break（释放） & 物体离开地面 & $\lambda_n\to0^+$ & 去激活约束行 \\
stick$\to$slip & 静摩擦被突破 & $\|\mathbf f_t\|\to\mu_f f_n$ & 约束变不等式 \\
slip$\to$stick & 滑动减速到零 & $\|\mathbf v_t\|\to0$ & 不等式变等式 \\
impact（碰撞） & 速度不连续 & $\mathbf v^+\neq\mathbf v^-$ & 冲量层求解 \\
\bottomrule
\end{tabular}
\end{table}

% ════════════════════════════════════════════════════════════════
\section{接触稳定性判据：CWC / ZMP / Capture Point}
\label{sec:cd-cwc}

\paragraph{从接触力到稳定性。}
有了接触约束的求解，腿足机器人还需判据：合接触力（wrench）是否落在「物理可行且不翻倒」的集合内。三大经典判据——ZMP、CWC、Capture Point——分别从「地反力矩零点」「接触 wrench 锥」「质心捕获」三个角度刻画。

\textbf{ZMP（Vukobratović--Stepanenko 1972）}\cite{vukobratovic1972stability}：地面上合地反力矩水平分量为零的点。以\emph{重力-惯性 wrench} 表述最为简洁——设地面参考点 $\mathbf p$、向上法向 $\mathbf n$，重力-惯性力 $\mathbf f^{gi}=m(\mathbf g-\ddot{\mathbf c})$ 与其对 $\mathbf p$ 的力矩 $\mathbf M_{\mathbf p}^{gi}=(\mathbf c-\mathbf p)\times m(\mathbf g-\ddot{\mathbf c})-\dot{\mathbf L}_G$（$\mathbf g$ 为指向下方的重力矢量，$\dot{\mathbf L}_G$ 为绕质心角动量变化率），则
\begin{equation}
  \mathbf p_{\text{ZMP}}=\mathbf p+\frac{\mathbf n\times\mathbf M_{\mathbf p}^{gi}}{\mathbf f^{gi}\cdot\mathbf n}.
  \label{eq:cd-zmp}
\end{equation}
分子的 $\mathbf n\times$ 把合力矩投影成地面内的位移、分母 $\mathbf f^{gi}\cdot\mathbf n$ 为法向反力（标量），故右端确为地面上一点的位置。\footnote{此处取重力-惯性 wrench 的标准写法（Sardain--Bessonnet 2004；亦见 Caron--Pham--Nakamura 的统一推导），与「合地反力矩水平分量为零」的原始 ZMP 定义等价、量纲一致——朴素地把合力矩本身当位移会缺少 $\mathbf n\times$ 投影与法向力归一。}

\begin{theorem}[Contact Wrench Cone，Caron--Pham--Nakamura 2015\cite{caron2015stability}]\label{thm:cd-cwc}
对 $2X\times2Y$ 矩形足，接触 wrench 可行性分解为三组：Coulomb $\|\mathbf f_t\|\le\mu_f f_z$；CoP 在支撑面内 $|\tau_x|\le Y f_z,\ |\tau_y|\le X f_z$；以及\emph{偏航（yaw）力矩界}——其紧界与切向力/力矩\emph{耦合}：
\begin{align}
  \tau_z^{\min}&=-\mu_f(X+Y)f_z+|Yf_x-\mu_f\tau_x|+|Xf_y-\mu_f\tau_y|,\nonumber\\
  \tau_z^{\max}&=+\mu_f(X+Y)f_z-|Yf_x+\mu_f\tau_x|-|Xf_y+\mu_f\tau_y|.
  \label{eq:cd-cwc-yaw}
\end{align}
\end{theorem}

\begin{note}
\textbf{yaw 界订正.}\;
$|\tau_z|\le\mu_f(X+Y)f_z$ 只是去掉 \cref{eq:cd-cwc-yaw} 耦合项后的\emph{最松（必要）界}，并非紧界——当耦合项 $|Yf_x\mp\mu_f\tau_x|,|Xf_y\mp\mu_f\tau_y|$ 为零时才退化到它。把它当作 tight 界会误导读者以为 yaw 与切向解耦。本书给出完整耦合式 \cref{eq:cd-cwc-yaw}。
\end{note}

\textbf{Capture Point（Pratt 2006）}\cite{pratt2006capture}：在线性倒立摆模型（LIPM）下使质心最终停到落足点的瞬时捕获点，
\begin{equation}
  \mathbf p_{\text{CP}}=\mathbf c+\dot{\mathbf c}/\omega,\qquad \omega=\sqrt{g/z_c},\qquad \text{DCM:}\ \dot{\boldsymbol\xi}=\omega(\boldsymbol\xi-\mathbf p_{\text{ZMP}}),\ \boldsymbol\xi=\mathbf c+\dot{\mathbf c}/\omega.
  \label{eq:cd-cp}
\end{equation}

\begin{pitfall}[编程陷阱：humanoid 足须 6D 面接触]
humanoid 足若当 3D 点接触，则无法约束 yaw 力矩（\cref{eq:cd-cwc-yaw} 退化），机器人会原地打转。正确做法是 6D 面接触；点接触仅适用于球面足。另：COP/ZMP/CP 三者在单足平地支撑下重合，但双足支撑或不平地面下分离——用 ZMP 做足力分配可能不稳。本判据主体接 \cref{ch:reduced-models}（LIPM/SRBD/CMM）与腿足部，互引。
\end{pitfall}

% ════════════════════════════════════════════════════════════════
\section{统一高地：Gauss 原理、Udwadia--Kalaba 与 IFT-on-KKT}
\label{sec:cd-gauss}

\paragraph{从特例回到原理。}
前面的 KKT、Baumgarte、闭链、接触都是「术」。本节登上「道」的高地：用一条 1829 年的变分原理统一解释为何约束加速度是「离自由加速度最近的」，由它导出无乘子的显式解，再用同一套 KKT 分解一举打通\emph{约束动力学的解析导数}——这正是把约束动力学焊到 MPC/DDP 上的关键接口。

\subsection{Gauss 最小约束原理与 Udwadia--Kalaba 方程}
\begin{theorem}[Gauss 最小约束原理 1829\cite{udwadia1992new}]\label{thm:cd-gauss}
受理想约束的真实加速度，在容许加速度集中\emph{最小化} Gauss 函数
\begin{equation}
  Z(\ddot{\mathbf q})=\tfrac12(\ddot{\mathbf q}-\mathbf f)^\top\mathbf M(\ddot{\mathbf q}-\mathbf f),\qquad \mathbf f=\mathbf M^{-1}(\boldsymbol\tau-\mathbf h),
  \label{eq:cd-gauss}
\end{equation}
即在 M-范数下约束加速度是「离自由加速度 $\mathbf f$ 最近的」。
\end{theorem}

\begin{theorem}[Udwadia--Kalaba 方程 1992\cite{udwadia1992new}]\label{thm:cd-uk}
给定线性约束 $\mathbf A\ddot{\mathbf q}=\mathbf b$，无需 Lagrange 乘子的显式解为
\begin{equation}
  \boxed{\;\ddot{\mathbf q}=\mathbf f+\mathbf M^{-1/2}(\mathbf A\mathbf M^{-1/2})^+(\mathbf b-\mathbf A\mathbf f)\;}
  \label{eq:cd-uk}
\end{equation}
其中 $(\cdot)^+$ 为 Moore--Penrose 伪逆。
\end{theorem}
\begin{proof}[骨架证明]
令 $\mathbf x=\mathbf M^{1/2}(\ddot{\mathbf q}-\mathbf f)$，约束变为 $\mathbf A\mathbf M^{-1/2}\mathbf x=\mathbf b-\mathbf A\mathbf f$；\cref{thm:cd-gauss} 即在此线性约束下最小化 $\|\mathbf x\|^2$，最小范数解 $\mathbf x=(\mathbf A\mathbf M^{-1/2})^+(\mathbf b-\mathbf A\mathbf f)$，代回得 \cref{eq:cd-uk}。\qed
\end{proof}

\begin{insight}[Gauss 是凸接触模型的合法性源头；WBC 是它的多任务推广]
Gauss 原理是 MuJoCo/Anitescu/Dojo 等凸接触模型的\emph{物理合法性源头}——它们都以某种形式最小化 Gauss 函数，只是约束集不同（等式 / 锥 / 互补）。而\textbf{WBC（整体控制）不是发明了新控制方法}：它是 Gauss 1829 思想在\emph{多任务、含不等式约束}场景下的现代推广——Gauss 在约束切空间中取「离自由加速度最近」的点，WBC 取「离期望任务加速度最近」的点，区别仅在距离度量（M-范数 vs 任务权重 $\mathbf W$）。把这条历史联系记牢，就不会把 WBC 当黑箱优化，而能从物理第一性原理理解它。
\end{insight}

\subsection{浮基方程与 WBC-QP 的数学骨架}
把上述统一视角落到腿足控制。设浮基坐标 $\mathbf q\in\SE(3)\times\mathbb R^n$、$n_v=n+6$，浮基动力学
\begin{equation}
  \mathbf M(\mathbf q)\ddot{\mathbf q}+\mathbf C\dot{\mathbf q}+\mathbf g=\mathbf S^\top\boldsymbol\tau+\mathbf J_c^\top\boldsymbol\lambda,\qquad \mathbf S=[\mathbf 0_{n\times6}\ \ \mathbf I_n],
  \label{eq:cd-floating}
\end{equation}
选择矩阵 $\mathbf S$ 剔除浮基 $6$ 行（浮基无电机）。\textbf{欠驱动}由此而来：浮基 $6$ 自由度无执行器，只能经接触力 $\mathbf J_c^\top\boldsymbol\lambda$ 间接控制；其前 $6$ 行恰是\emph{质心动力学}（Orin--Goswami--Lee 2013\cite{orin2013centroidal}）$\dot{\mathbf h}_G=[m\ddot{\mathbf c};\dot{\mathbf L}_G]$，质心动量 $\mathbf h_G=\mathbf A_G(\mathbf q)\mathbf v$（CMM，$\mathbf A_G\in\mathbb R^{6\times n_v}$ 为整机惯量前 $6$ 行经质心坐标变换，cref \cref{ch:reduced-models}）。WBC-QP 以 $(\ddot{\mathbf q},\boldsymbol\tau,\boldsymbol\lambda)$ 为决策变量：
\begin{equation}
  \min_{\ddot{\mathbf q},\boldsymbol\tau,\boldsymbol\lambda}\ \sum_k w_k\big\|\mathbf J_k\ddot{\mathbf q}+\dot{\mathbf J}_k\dot{\mathbf q}-\ddot{\mathbf x}_k^{\text{des}}\big\|^2\quad\text{s.t.}\quad \text{\cref{eq:cd-floating}},\ \mathbf J_c\ddot{\mathbf q}+\dot{\mathbf J}_c\dot{\mathbf q}=\mathbf 0,\ \boldsymbol\lambda\in\mathcal{FC}(\mu_f),\ \boldsymbol\tau\in[\boldsymbol\tau_{\min},\boldsymbol\tau_{\max}].
  \label{eq:cd-wbc}
\end{equation}
由 \cref{eq:cd-floating} 可消去 $\boldsymbol\tau=(\mathbf S\mathbf S^\top)^{-1}\mathbf S(\mathbf M\ddot{\mathbf q}+\mathbf C\dot{\mathbf q}+\mathbf g-\mathbf J_c^\top\boldsymbol\lambda)$，决策变量降到 $(\ddot{\mathbf q},\boldsymbol\lambda)$。严格优先级用分层 QP（HQP）实现：级联（先解高优先、再约束其不变坏）/ 零空间投影（低优先在高优先零空间内优化）/ 软权重（$w_1\gg w_2\gg\cdots$，简单但不严格）。接触力分配的冗余（四足 $4\times6=24$ 维 $\boldsymbol\lambda$ 对 $6$ 维质心方程）由 Gauss 型准则（最小力范数等）选优。\textbf{WBC 库实现主体（TSID/HQP 接口、QP 求解器选型）在 \cref{ch:force-wbc}}，本章只给数学骨架与 Gauss 最优性。

\subsection{IFT-on-KKT：约束动力学的解析导数}
最后，同一套 KKT 分解还能\emph{免费}给出解析导数。设 KKT 残差 $\mathbf F(\mathbf z,\boldsymbol\theta)=\mathbf 0$，$\mathbf z=(\ddot{\mathbf q},\boldsymbol\lambda)$，$\boldsymbol\theta\in\{\mathbf q,\mathbf v,\boldsymbol\tau\}$：
\begin{equation}
  \mathbf F=\begin{bmatrix}\mathbf M\ddot{\mathbf q}+\mathbf h-\mathbf S^\top\boldsymbol\tau-\mathbf J^\top\boldsymbol\lambda\\ \mathbf J\ddot{\mathbf q}+\boldsymbol\gamma\end{bmatrix}=\mathbf 0\ \Longrightarrow\ \frac{\partial\mathbf z}{\partial\boldsymbol\theta}=-\underbrace{\begin{bmatrix}\mathbf M&-\mathbf J^\top\\\mathbf J&\mathbf 0\end{bmatrix}^{-1}}_{\text{forward 已分解}}\frac{\partial\mathbf F}{\partial\boldsymbol\theta}.
  \label{eq:cd-ift-kkt}
\end{equation}

\begin{insight}[KKT 分解的复用]
\cref{eq:cd-ift-kkt} 的妙处：KKT 矩阵的分解\emph{既}用于 forward（解出 $\ddot{\mathbf q},\boldsymbol\lambda$）\emph{又}用于 derivative——每个参数方向的梯度只需一次回代 $O((n_v+m)^2)$，无需重新分解。这是「解一次、求导也几乎免费」的范式，正是 \cref{ch:ddp-ilqr} DDP 实时性的来源。
\end{insight}

\begin{theorem}[Pinocchio 3 proximal 天然可微 2021\cite{carpentier2021proximal}]\label{thm:cd-proximal}
把 KKT 正则化为 $\begin{psmallmatrix}\mathbf M&-\mathbf J^\top\\\mathbf J&\mu\mathbf I\end{psmallmatrix}$ 并配增广 Lagrange 迭代 $\boldsymbol\lambda^{(k+1)}=\boldsymbol\lambda^{(k)}+\tfrac1\mu(\mathbf J\ddot{\mathbf q}^{(k+1)}+\boldsymbol\gamma)$，则 $(\mathbf G+\mu\mathbf I)$ 对任意 $\mathbf J$（即便秩亏）\emph{恒可逆}，故 IFT \cref{eq:cd-ift-kkt} 永远成立；且 $\mu\to0$ 时 $(\mathbf G+\mu\mathbf I)^{-1}\to\mathbf G^+$，正则化伪逆收敛到 Udwadia--Kalaba 的伪逆 \cref{eq:cd-uk}。
\end{theorem}
\noindent 这把 \cref{thm:cd-uk}（UK 显式解）、\cref{eq:cd-ift-kkt}（IFT 可微）、LICQ 鲁棒性三者统一在一个正则化参数 $\mu$ 下。\textbf{务必区分} proximal 的 $\mu$ 与摩擦系数 $\mu_f$：前者是数值正则化（正常 $10^{-10}$–$10^{-8}$、秩亏时增到 $10^{-4}$；太大引入伪滑移、太小数值不稳），后者是物理量。

% ════════════════════════════════════════════════════════════════
\section{可微接触的四条路线}
\label{sec:cd-diffcontact}

\paragraph{动机：端到端可微的最后断点。}
2020 年以来机器人学的一大趋势是\emph{端到端可微}——从感知到控制到仿真整条链都能反传梯度。但接触是这条链上最大的「断点」：LCP 的互补条件 $\lambda\phi=0$ 在切换点（接触建立/释放、stick$\leftrightarrow$slip）处\emph{不可微}。若接触不可微，policy gradient 过仿真器的梯度为零或不存在、model-based RL 学不到接触瞬间该怎么做、轨迹优化在接触切换处 NLP 梯度断裂。四条路线给出不同方案，各有代价（\cref{tab:cd-diff}）。

\begin{itemize}
  \item \textbf{(A) MuJoCo 软化}\cite{todorov2014convex}：放弃严格互补，用 $\mathbf R>\mathbf 0$ 正则化把问题变成凸 QP/SOCP $\min\tfrac12\mathbf x^\top\mathbf M\mathbf x-\mathbf x^\top\boldsymbol\tau_{\text{bias}}+\tfrac12\|\mathbf y-\mathbf a^*\|_{\mathbf R}^2\ \text{s.t.}\ \mathbf y=\mathbf J\mathbf x+\mathbf c,\ \mathbf y\in\mathcal K^*$。凸、解析可逆、梯度良定（MJX 原生支持），但物理不精确（允许轻微穿透与人工滑移）。
  \item \textbf{(B) Dojo $\kappa$-平滑}\cite{howell2022dojo}：保持硬接触物理，用内点法参数 $\kappa$ 平滑互补 $\mathbf s\circ\boldsymbol\lambda=\kappa\mathbf e$（$\kappa\to0$ 回到精确互补）；对 KKT 残差用 IFT 求梯度。$\kappa$ 是\emph{用户可调的精度-梯度质量权衡}：大 $\kappa$ 平滑、梯度信息丰富但物理不够精确，小 $\kappa$ 精确但切换点附近梯度陡。
  \item \textbf{(C) Nimble 解析子梯度}\cite{werling2021fast}：保留 hard-LCP、不做任何平滑。在 active set 不切换的区间内，接触力是参数的\emph{解析}函数（LCP 退化为线性方程组、求导平凡）；在切换点用\emph{子梯度}（取左导或右导之一）。物理完全精确，但切换点处梯度可能不连续，需后续插值或 bundle 方法。
  \item \textbf{(D) 随机平滑}\cite{suh2022differentiable}：不修改仿真器，而对输入参数加随机扰动取期望 $\nabla_{\boldsymbol\theta}\mathbb E_{\boldsymbol\epsilon}[L(\text{sim}(\boldsymbol\theta+\boldsymbol\epsilon))]$，即 score-function/REINFORCE 估计——不需仿真器可微，只需能采样，代价是高方差。
\end{itemize}

\begin{table}[H]\centering\small
\caption{可微接触四路线对比}
\label{tab:cd-diff}
\begin{tabularx}{\linewidth}{p{0.16\linewidth}LLLL}
\toprule
维度 & MuJoCo 软化 & Dojo $\kappa$-smooth & Nimble 子梯度 & 随机平滑 \\
\midrule
物理精度 & 中（有穿透/滑移） & 高（$\kappa\to0$ 精确） & 高（精确 LCP） & 依赖底层 sim \\
梯度质量 & 好（凸解析） & 好（可调 $\kappa$） & 切换点不连续 & 高方差 \\
求解代价 & 低（凸 QP） & 中（IPM） & 中（稀疏线性） & 高（多次采样） \\
适用场景 & GPU 大批量 RL & MPC/TO & 精确物理需求 & 任意不可微 sim \\
\bottomrule
\end{tabularx}
\end{table}

\begin{pitfall}[思维陷阱：「可微仿真 = 更好的仿真」]
可微化有代价（平滑化降精度、子梯度有噪声）。对不需梯度的任务（如 RL 用 reward shaping），高速仿真的速度优势可能更重要。可微仿真是\emph{工具}不是\emph{目标}，按任务选方法。注意 Dojo 的 $\kappa$ 与 Pinocchio 的 $\mu$ 在数学上等价——都是互补条件的正则化。Le~Lidec--Jallet--Carpentier（TRO 2024）\cite{lelidec2024contact} 对主流接触模型做了系统比较，可参。
\end{pitfall}

% ════════════════════════════════════════════════════════════════
\section{动力学解析微分：需求与四条路线}
\label{sec:cd-deriv-need}

\paragraph{承上：从「求解约束」到「对约束求导」。}
\cref{sec:cd-gauss} 的 IFT-on-KKT 已点明：约束动力学的解析导数是把它焊到最优控制的关键。但 IFT 需要 $\partial\mathbf F/\partial\boldsymbol\theta$，其中藏着对\emph{无约束}动力学（RNEA/ABA）求导这块硬骨头。本部分（\cref{sec:cd-deriv-need}--\cref{sec:cd-second}）专攻它：先讲为何动力学导数是 MPC 的瓶颈、四条计算路线如何取舍，再给 $\partial\mathrm{RNEA}/\partial\mathbf q$ 的 $O(Nd)$ 解析递推与 $\partial\mathrm{ABA}$ 的 ID--FD 互逆。

\paragraph{动机：动力学导数是 MPC 的瓶颈。}
DDP/iLQR 的 Riccati 反传需要在每个时间节点求 $\partial\mathbf f/\partial\mathbf x$ 与 $\partial\mathbf f/\partial\mathbf u$（\cref{ch:ddp-ilqr}）。若用最朴素的有限差分：对 $36$ 自由度人形，每次中心差分 $\partial\boldsymbol\tau/\partial\mathbf q$ 需 $73$ 次 RNEA。乘以时域节点数（$50$–$100$）、iLQR 迭代数（$5$–$20$）、MPC 频率（$100$\,Hz），每秒就要数千到上万次导数评估——有限差分路线让 CPU 占用爆到几百个百分点（完全不可能实时），而解析路线把同样的工作压到微秒级、占用降到几十个百分点。这不是「快一点」，而是\emph{实时可行 vs 完全不可行}的分水岭。动力学导数还支撑轨迹优化、系统辨识（Fisher 信息 $\mathbf Y^\top\mathbf Y$）、co-design、model-based RL、可微渲染等。

\begin{insight}[解析导数严格保零模式，对 DDP 稀疏性至关重要]
当系统静止（$\dot{\mathbf q}=\ddot{\mathbf q}=\mathbf 0$）且无重力时，$\partial\boldsymbol\tau/\partial\dot{\mathbf q}=\mathbf C(\mathbf q,\mathbf 0)=\mathbf 0$ 理论上\emph{严格}为零。有限差分却产生约 $10^{-7}$ 的\emph{稠密残差}：DDP 的 Riccati 反传把这些「伪非零」当真，误判质量矩阵的稀疏结构（把本应带状的 $O(Nd)$ 当稠密 $O(N^2)$），QP 求解器速度直接掉一个数量级，控制器在平衡点附近持续微抖。解析导数则精确返回零矩阵——因为它直接从递推公式的代数结构计算，输入为零时相关项自动消失。
\end{insight}

\subsection{四条路线的取舍}
\textbf{(1) 有限差分}：$\partial f/\partial x_j\approx[f(x+\varepsilon e_j)-f(x)]/\varepsilon$（前向）或中心差分。误差由\emph{截断}（$\sim\tfrac\varepsilon2 f''$）与\emph{舍入}（$\sim\varepsilon_{\text{mach}}/\varepsilon$）两部分构成，最优步长前向 $\varepsilon^*=\sqrt{\varepsilon_{\text{mach}}}\approx10^{-8}$、中心 $\varepsilon^*=\varepsilon_{\text{mach}}^{1/3}\approx5\times10^{-6}$；前向需 $(n{+}1)$ 次、中心需 $2n$ 次 RNEA。\emph{致命缺陷}：破坏零模式（见上）。验证解析导数时不可替代。\;
\textbf{(2) 运行时自动微分（AD）}：基于\emph{对偶数} $a+b\varepsilon\ (\varepsilon^2=0)$，因 $f(a+b\varepsilon)=f(a)+f'(a)b\varepsilon$（Taylor 展开高阶项因 $\varepsilon^2=0$ 全灭），设种子 $b=1$ 即得 $f'(a)$，机器精度、无步长问题。Ceres 的 Jet 把 $\varepsilon$ 推广为 $N$ 维同时追踪 $N$ 个方向；CppAD 用算子重载录制计算图（tape）回放。开销约原生 RNEA 的 $5$–$10$ 倍。\;
\textbf{(3) 代码生成（CodeGen）}：把 tape「打印」成纯 C、编译为 \texttt{.so}，消除所有乘 $0$/加 $0$/乘 $1$、稀疏 Jacobian 只生成非零项。机器精度、接近手写 C，但有首次编译开销（数秒到数分钟）。\;
\textbf{(4) 解析微分}：直接对递推公式求导。\emph{零编译开销 $+$ 即时可用 $+$ 机器精度 $+$ 严格保零模式}——下两节的主题。\cref{tab:cd-deriv} 定性对比（所有绝对耗时随 CPU/版本/编译选项浮动，此处只取定性量级）。

\begin{table}[H]\centering\small
\caption{四条导数路线的定性对比}
\label{tab:cd-deriv}
\begin{tabular}{p{0.20\linewidth}p{0.22\linewidth}p{0.18\linewidth}p{0.14\linewidth}p{0.12\linewidth}}
\toprule
路线 & 相对耗时（量级） & 首次开销 & 精度 & 保零模式 \\
\midrule
有限差分 & 慢（$\sim2n\times$ RNEA） & 无 & $\sim10^{-7}$ & 否 \\
运行时 AD & 中（$5$–$10\times$ RNEA） & 无 & 机器级 & 是 \\
CodeGen & 快（$\sim$ 手写 C） & 数秒–分钟 & 机器级 & 是 \\
解析微分 & 快（$\sim3\times$ RNEA） & 无 & 机器级 & 是 \\
\bottomrule
\end{tabular}
\end{table}

\begin{note}
\textbf{范围说明.}\;
Pinocchio/CppAD/CasADi/JAX 的具体 API、CppADCodeGen 六步代码、嵌入式部署、性能基准 $\mu$s 表属工程实现，本书略去，仅一句点名「该框架的工业实现」。其中唯一可精确引用的绝对数是 Todorov 2014\cite{todorov2014convex}：$27$ 自由度 humanoid $+\ 10$ 接触正向求解约 $0.1$\,ms（$100\times$ 实时），原论文有载。
\end{note}

% ════════════════════════════════════════════════════════════════
\section{\texorpdfstring{$\partial\mathrm{RNEA}/\partial\mathbf q$}{dRNEA/dq} 的 O(Nd) 闭式}
\label{sec:cd-drnea}

\paragraph{技术核心。}
Carpentier--Mansard（RSS 2018）\cite{carpentier2018analytical} 给出 RNEA 一阶偏导的 $O(Nd)$ 解析递推。其「一招制胜」的代数技巧是：\textbf{placement（Plücker 变换）对关节角的微分 $=$ 被作用的空间量与运动子空间 $\mathbf S_i$ 的空间叉乘}——这条恒等式把看似需要 $6\times6$ 张量求导的问题压成一次熟悉的空间叉乘。

\subsection{RNEA 回顾与 placement 微分恒等式}
回顾 \cref{ch:spatial-dynamics} 的 RNEA（Featherstone 记号，$\lambda(i)$ 为 $i$ 的父关节）。前向 pass（$i=1..N$）：$\mathbf v_i={}^i\mathbf X_{\lambda(i)}\mathbf v_{\lambda(i)}+\mathbf S_i\dot q_i$、$\mathbf a_i={}^i\mathbf X_{\lambda(i)}\mathbf a_{\lambda(i)}+\mathbf S_i\ddot q_i+\mathbf v_i\times(\mathbf S_i\dot q_i)$、$\mathbf h_i=\mathcal I_i\mathbf v_i$、$\mathbf f_i=\mathcal I_i\mathbf a_i+\mathbf v_i\times^*\mathbf h_i-\mathbf f_i^{\text{ext}}$；后向 pass（$i=N..1$）：$\tau_i=\mathbf S_i^\top\mathbf f_i$、$\mathbf f_{\lambda(i)}\mathrel{+}={}^{\lambda(i)}\mathbf X_i^*\mathbf f_i$。

\begin{theorem}[placement 微分恒等式，Carpentier--Mansard 2018 Eq.13--14\cite{carpentier2018analytical}]\label{thm:cd-placement}
\begin{equation}
  \frac{\partial({}^i\mathbf X_{\lambda(i)}\mathbf v_{\lambda(i)})}{\partial q_i}=\big({}^i\mathbf X_{\lambda(i)}\mathbf v_{\lambda(i)}\big)\times\mathbf S_i,\qquad
  \frac{\partial({}^{\lambda(i)}\mathbf X_i^*\mathbf f_i)}{\partial q_i}={}^{\lambda(i)}\mathbf X_i^*\,(\mathbf S_i\times^*\mathbf f_i).
  \label{eq:cd-placement}
\end{equation}
\end{theorem}
\begin{proof}[骨架证明（revolute 关节）]
设 ${}^i\mathbf X_{\lambda(i)}=\mathbf X_{\text{fixed}}\cdot\mathbf X_{\text{joint}}(q_i)$，$\mathbf X_{\text{joint}}$ 只含绕旋转轴 $\mathbf s$ 的旋转。对 $q_i$ 求导，$\mathrm d\mathbf X_{\text{joint}}/\mathrm dq_i$ 的作用等价于在 child frame 做无穷小旋转 $\hat{\mathbf s}$——这正是 $\mathbf S_i$ 的定义。故关节 $i$ 转 $\mathrm dq_i$ 时连杆 $i$ 的 frame 微旋，其对任何空间量的效应 $=$ 将该量与旋转方向 $\mathbf S_i$ 做一次空间叉乘。\qed
\end{proof}

\begin{insight}[一招制胜]
\cref{eq:cd-placement} 是整篇论文的核心：把「$6\times6$ 矩阵 ${}^i\mathbf X$ 关于标量 $q_i$ 的导数」这个看似要张量运算的对象，压缩成一次空间叉乘 $\times/\times^*$（即 $\mathrm{ad}/\mathrm{ad}^*$，见 \cref{ch:lie}）。整个 $O(Nd)$ 复杂度由此而来。
\end{insight}

\begin{note}
\textbf{右扰动对齐（J$\leftrightarrow$L 跨章最易错处）.}\;
\cref{eq:cd-placement} 用 Featherstone $[\boldsymbol\omega;\mathbf v]$ 序；但本书\emph{对 $\mathbf q$ 的导数取右扰动} $\mathbf R\,\mathrm{Exp}(\delta\boldsymbol\phi)$（右雅可比 $\mathbf J_r$，与 \cref{ch:lie} 一致）。叉乘左右顺序、符号、${}^i\mathbf X$ 变换方向依赖空间向量排序与扰动约定；本书 $\se(3)$ 序 $\boldsymbol\xi=[\boldsymbol\rho;\boldsymbol\phi]$ 与 $[\boldsymbol\omega;\mathbf v]$ 由置换 $\mathbf P=\begin{psmallmatrix}\mathbf 0&\mathbf I\\\mathbf I&\mathbf 0\end{psmallmatrix}$ 互换（ad 块的上下三角随之互换）。延续 \cref{ch:spatial-dynamics} 已建的 $[\boldsymbol\omega;\mathbf v]\leftrightarrow\boldsymbol\xi=[\boldsymbol\rho;\boldsymbol\phi]$ 对照，每个解析导数公式中对 $\mathbf q$ 的求导均按右扰动理解。浮基时这些导数实际在\emph{切空间}上表达（$\partial\boldsymbol\tau/\partial\mathbf q|_{\text{tangent}}$，对切空间变量 $\mathbf v$ 而非参数 $\mathbf q$ 求），返回 $n_v\times n_v$ 矩阵；用有限差分验证时须经 Lie 群积分 $\mathrm{Exp}(\mathbf q,h\mathbf e_j)$ 扰动，直接 $\mathbf q+h\mathbf e_j$ 会破坏四元数单位约束。
\end{note}

\subsection{前向与后向导数 pass}
对微分变量 $u\in\{\mathbf q,\dot{\mathbf q}\}$，\textbf{前向 pass} 沿树递推：速度导数 $\partial\mathbf v_i/\partial u={}^i\mathbf X_{\lambda(i)}\partial\mathbf v_{\lambda(i)}/\partial u+(\text{placement 项})$，其中对 $u=q_i$（自身关节）由 \cref{eq:cd-placement} 添加 $({}^i\mathbf X_{\lambda(i)}\mathbf v_{\lambda(i)})\times\mathbf S_i$，对 $u=\dot q_i$ 直接添加 $\mathbf S_i$；动量导数 $\partial\mathbf h_i/\partial u=\mathcal I_i\partial\mathbf v_i/\partial u$；加速度与力导数同理递推（含 $\partial\mathbf v_i/\partial u\times\mathbf v_J$、$\partial\mathbf v_i/\partial u\times^*\mathbf h_i$ 等交叉项）。\textbf{后向 pass}：$\partial\tau_i/\partial u=\mathbf S_i^\top\partial\mathbf f_i/\partial u$（填到 $\partial\boldsymbol\tau/\partial u$ 第 $i$ 行），并向父节点传播 $\partial\mathbf f_{\lambda(i)}/\partial u\mathrel{+}={}^{\lambda(i)}\mathbf X_i^*\partial\mathbf f_i/\partial u$，对 $u=q_i$ 额外加 dual placement 项 ${}^{\lambda(i)}\mathbf X_i^*(\mathbf S_i\times^*\mathbf f_i)$。下框是语言无关算法骨架。

\begin{algo}[$\partial\mathrm{RNEA}/\partial(\mathbf q,\dot{\mathbf q})$ 前向+后向 pass\cite{carpentier2018analytical}]\label{alg:cd-drnea}
\begin{flushleft}\small\ttfamily
// 前向 pass（i = 1..N，X = iX\_lambda(i)，S = S\_i）\\
\quad dv\_i/dq = X · dv\_lambda(i)/dq;\ \ dv\_i/dq[:,i] += (X·v\_lambda(i)) × S\\
\quad dv\_i/dq̇ = X · dv\_lambda(i)/dq̇;\ \ dv\_i/dq̇[:,i] += S\\
\quad da\_i/dq = X · da\_lambda(i)/dq + (X·a\_lambda(i))×S [列 i] + dv\_i/dq × v\_J + v\_i × dv\_J/dq\\
\quad dh\_i/du = I\_i · dv\_i/du;\ \ df\_i/du = I\_i·da\_i/du + dv\_i/du ×* h\_i + v\_i ×* dh\_i/du\\
// 后向 pass（i = N..1）\\
\quad dtau\_i/du = S\_i\textasciicircum T · df\_i/du\ \ // 第 i 行\\
\quad if lambda(i)!=0:\ \ df\_lambda(i)/du += X*·df\_i/du;\ \ df\_lambda(i)/dq[:,i] += X*·(S\_i ×* f\_i)
\end{flushleft}
\end{algo}

\subsection{三个关键简化}
\textbf{简化 1：$\partial\boldsymbol\tau/\partial\dot{\mathbf q}$ 比 $\partial\boldsymbol\tau/\partial\mathbf q$ 更简。} placement 不依赖速度（$\partial\mathbf X_J/\partial\dot q=0$），$\partial\mathbf v_J/\partial\dot q_i=\mathbf S_i$（线性），故加速度方程中许多项消失——velocity 对 $\dot{\mathbf q}$ 线性、二阶导自动为零。\;
\textbf{简化 2：$\partial\boldsymbol\tau/\partial\ddot{\mathbf q}=\mathbf M(\mathbf q)$。} 由 $\boldsymbol\tau=\mathbf M\ddot{\mathbf q}+\mathbf C\dot{\mathbf q}+\mathbf g$ 直接得，等于 CRBA 输出的质量矩阵，无需额外推导。\;
\textbf{简化 3：稀疏性\emph{仅对} $\partial\boldsymbol\tau/\partial\ddot{\mathbf q}=\mathbf M$ 严格成立。}

\begin{pitfall}[概念误区订正：$\partial\boldsymbol\tau/\partial\mathbf q$ 串联链稠密，$O(Nd)$ 是算法复杂度而非结果稀疏]
一个流行的旧误解是「$\partial\tau_i/\partial q_j$ 只在 $j$ 是 $i$ 祖先时非零，故串联臂 $\partial\boldsymbol\tau/\partial\mathbf q$ 下三角」。\emph{这是错的}。反例：2R 平面臂 $\tau_1=M_{11}\ddot q_1+M_{12}\ddot q_2+C_{112}\dot q_1\dot q_2+C_{122}\dot q_2^2+g_1(q_1,q_2)$，其 $\partial\tau_1/\partial q_2\neq0$（经重力与 Coriolis 项），尽管关节 $2$ 是关节 $1$ 的\emph{后代}。正确陈述：质量矩阵 $\mathbf M(\mathbf q)=\partial\boldsymbol\tau/\partial\ddot{\mathbf q}$ 由树拓扑结构稀疏（分支树 $M_{ij}=0\iff i,j$ 不在同一根-叶路径），这是\emph{真正}的结构稀疏；$\partial\boldsymbol\tau/\partial\mathbf q,\partial\boldsymbol\tau/\partial\dot{\mathbf q}$ 对串联链（如 KUKA $7\times7$）\emph{通常稠密}；$O(Nd)$ 复杂度源于\emph{算法}按列只沿关节到根的深度-$d$ 路径传播，\emph{不}依赖结果矩阵稀疏。
\end{pitfall}

% ════════════════════════════════════════════════════════════════
\section{\texorpdfstring{$\partial\mathrm{ABA}$}{dABA} via ID--FD 互逆与 AZA}
\label{sec:cd-daba}

\paragraph{动机：避开三阶张量。}
正向动力学 $\ddot{\mathbf q}=\mathrm{FD}(\mathbf q,\mathbf v,\boldsymbol\tau)$，DDP/iLQR 需 $\partial\ddot{\mathbf q}/\partial(\mathbf q,\mathbf v,\boldsymbol\tau)$。直接对 ABA 三趟递推求导会引入大量中间张量（铰接体惯量 $\mathcal I_i^A$、偏置力 $\mathbf p_i^A$ 本身是 $6\times6$ 矩阵，求导产生 $6\times6\times n$ 三阶张量）。Carpentier--Mansard 的妙招是利用 ID 与 FD 的\emph{互逆关系}。

\begin{theorem}[ID--FD 互逆\cite{carpentier2018analytical}]\label{thm:cd-idfd}
逆动力学 ID 与正动力学 FD 互逆：$\mathrm{ID}(\mathbf q,\mathbf v,\mathrm{FD}(\mathbf q,\mathbf v,\boldsymbol\tau))=\boldsymbol\tau$。对 $u\in\{\mathbf q,\mathbf v\}$ 求偏导（$\boldsymbol\tau$ 不变），并用 $\partial\mathrm{ID}/\partial\ddot{\mathbf q}=\partial\boldsymbol\tau/\partial\ddot{\mathbf q}=\mathbf M(\mathbf q)$，得
\begin{equation}
  \boxed{\;\frac{\partial\ddot{\mathbf q}}{\partial\mathbf q}=-\mathbf M^{-1}\frac{\partial\boldsymbol\tau}{\partial\mathbf q},\qquad \frac{\partial\ddot{\mathbf q}}{\partial\mathbf v}=-\mathbf M^{-1}\frac{\partial\boldsymbol\tau}{\partial\mathbf v},\qquad \frac{\partial\ddot{\mathbf q}}{\partial\boldsymbol\tau}=\mathbf M^{-1}.\;}
  \label{eq:cd-idfd}
\end{equation}
\end{theorem}

\begin{insight}[正向导数化为逆向导数 + 线性求解]
\cref{eq:cd-idfd} 把 ABA 导数问题\emph{完全还原}为 RNEA 导数（\cref{sec:cd-drnea} 已算好）加一次 $\mathbf M^{-1}$ 矩阵乘——避免了直接微分 ABA 三趟递推产生的三阶张量。这正是「把正向导数问题转化为逆向导数 $+$ 线性求解」的范式。
\end{insight}

\subsection{高效 \texorpdfstring{$\mathbf M^{-1}$}{Minv}：稀疏分解与 AZA}
\cref{eq:cd-idfd} 需要高效的 $\mathbf M^{-1}$。两条路：\textbf{(1) CRBA $+$ 稀疏 $\mathbf L\mathbf D\mathbf L^\top$ 分解}（利用运动树的消元顺序，$O(Nd^2)$，与 \cref{ch:isam2} SLAM 后端的稀疏 Cholesky 同源）；\textbf{(2) AZA 算法}——ABA 在 $\boldsymbol\tau=\mathbf e_j$、$\dot{\mathbf q}=\mathbf 0$、$\mathbf g=\mathbf 0$ 的特殊输入下，输出恰为 $\mathbf M^{-1}$ 的第 $j$ 列：
\begin{equation}
  \ddot{\mathbf q}=\mathrm{ABA}(\mathbf q,\mathbf 0,\mathbf e_j)=\mathbf M^{-1}\mathbf e_j=(\mathbf M^{-1})_{:,j},
  \label{eq:cd-aza}
\end{equation}
对 $j=1..n$ 逐列、利用对称只算上三角即得完整 $\mathbf M^{-1}$。

\begin{note}
\textbf{AZA 归属订正.}\;
\cref{eq:cd-aza} 这一 $\mathbf M^{-1}$ 直接算法的解析形式\emph{首见} Carpentier 2018（\emph{Analytical Inverse of the JSIM}）\cite{carpentier2018inverse}；Singh--Russell--Wensing 2021（\emph{Efficient Analytical Derivatives of RBD using Spatial Vector Algebra}，arXiv:2105.05102）\cite{singh2021efficient} 用空间向量代数整理并命名为 \textbf{AZA}（ABA-Zero-Algorithm）。AZA 因 ABA 递推天然利用树结构、省去 Cholesky 的填入处理，通常比 CRBA$+$Cholesky 快约一倍。
\end{note}

\subsection{装配 DDP 的 Jacobian}
DDP 连续时间状态方程 $\dot{\mathbf x}=\mathbf f(\mathbf x,\mathbf u)=[\mathbf v;\,\mathrm{FD}(\mathbf q,\mathbf v,\boldsymbol\tau)]$，其 Jacobian
\begin{equation}
  \mathbf A=\frac{\partial\mathbf f}{\partial\mathbf x}=\begin{bmatrix}\mathbf 0 & \mathbf I\\ \partial\ddot{\mathbf q}/\partial\mathbf q & \partial\ddot{\mathbf q}/\partial\mathbf v\end{bmatrix},\qquad \mathbf B=\frac{\partial\mathbf f}{\partial\mathbf u}=\begin{bmatrix}\mathbf 0\\ \mathbf M^{-1}\end{bmatrix},
  \label{eq:cd-ddp-jac}
\end{equation}
恰好就是 \cref{eq:cd-idfd} 的三个输出（浮基时 $\mathbf M^{-1}$ 为 $(6{+}n)\times(6{+}n)$，而 $\mathbf B$ 只取被驱动列、前 $6$ 行对应欠驱动浮基）。\textbf{这些 Jacobian 的最优控制应用——iLQR/DDP 的 Riccati 反传——见 \cref{ch:ddp-ilqr}}；本章记号 $\mathbf A,\mathbf B$ 与 \cref{ch:ddp-ilqr} 的 $\mathbf f_x,\mathbf f_u$ 一致。

% ════════════════════════════════════════════════════════════════
\section{二阶导数、Exact DDP 与惯性辨识}
\label{sec:cd-second}

\subsection{二阶导数与 Exact DDP}
iLQR 用 Gauss--Newton 近似 $\mathbf Q_{xx}\approx\boldsymbol\ell_{xx}+\mathbf f_x^\top\mathbf V_{xx}'\mathbf f_x$，\emph{弃掉} $\mathbf V_x'\cdot\mathbf f_{xx}$ 项；\textbf{exact DDP} 保留 $\sum_k(\mathbf V_x')_k(\mathbf f_{xx})_k$，获得二次收敛。何时需要 exact：大扰动、非最小二乘 cost、约束增广 Lagrange。

\begin{insight}[GN 在解处等于 exact]
在最优解处伴随 $\mathbf V_x'=\mathbf 0$，被弃项自动消失，故 Gauss--Newton 在收敛点与 exact DDP 一致——这解释了为何 iLQR 在轨迹收敛后精度不输 exact DDP，而在远离解处 exact 的二次收敛更快。
\end{insight}

二阶项需要动力学的二阶偏导 $\partial^2\mathrm{ID}/\partial\mathbf q^2$ 等。Singh--Russell--Wensing（IROS 2022，arXiv:2203.01497）\cite{singh2022second} 给出逆动力学二阶偏导的 $O(Nd^2)$ 解析递推，把 exact DDP 的二阶张量计算也压到与一阶同量级。

\subsection{惯性辨识的 Regressor}
动力学导数的另一大应用是系统辨识。Lagrange 方程对惯性参数 $\boldsymbol\theta=(m_i,m_i\mathbf c_i,\mathbf I_i)$ \emph{线性}：
\begin{equation}
  \boldsymbol\tau=\mathbf Y(\mathbf q,\dot{\mathbf q},\ddot{\mathbf q})\,\boldsymbol\theta,\qquad \mathbf Y\in\mathbb R^{n_v\times10n},\ \boldsymbol\theta\in\mathbb R^{10n},
  \label{eq:cd-regressor}
\end{equation}
其中每连杆 $10$ 个参数（质量 $1$ $+$ 质心 $3$ $+$ 惯性张量上三角 $6$），$\mathbf Y$ 为回归矩阵。线性性源于 RNEA 力递推 $\mathbf f_i=\mathcal I_i\mathbf a_i+\mathbf v_i\times^*(\mathcal I_i\mathbf v_i)-\mathbf f_i^{\text{ext}}$：空间惯量 $\mathcal I_i$ 对 $10$ 参数线性，而 $\mathbf a_i,\mathbf v_i$ 不依赖惯性参数，故 $\mathbf f_i$ 进而 $\tau_i=\mathbf S_i^\top\mathbf f_i$ 对 $\boldsymbol\theta$ 线性。$10n$ 个参数中许多不可辨识，\textbf{base parameter reduction}（Gautier--Khalil QR 分解）\cite{khalil2002modeling} 给出通常约 $4$–$6n$ 个基参数。辨识精度要求\emph{持续激励} $\int_0^T\mathbf Y^\top\mathbf Y\,\mathrm dt\succ\mathbf 0$（需设计 Fourier 级数等激励轨迹）；Fisher 信息 $\mathcal I=\sum_t\mathbf Y_t^\top\boldsymbol\Sigma^{-1}\mathbf Y_t$ 与 Cramér--Rao 下界 $\mathcal I^{-1}$ 把它与 \cref{ch:isam2} 估计部的信息矩阵焊接。运动学导数（$\dot{\mathbf J}$、运动学 Hessian）在操作空间控制中的用途主体留 \cref{ch:operational-space}。

% ════════════════════════════════════════════════════════════════
\section{本章脉络与展望}
\label{sec:cd-outlook}

本章把约束动力学从 18 世纪分析力学一路打通到 21 世纪可微仿真与 MPC 梯度。主线如下：

\begin{center}
\small
\begin{tikzpicture}[node distance=2.6mm and 5mm, >=Stealth,
  b/.style={draw, rounded corners, align=center, font=\scriptsize, inner sep=2.5pt, fill=main!6, draw=main!60!black}, ln/.style={->, thick, gray!55!black}]
\node[b] (a) {约束分类};
\node[b, right=of a] (c) {DAE/指数};
\node[b, right=of c] (d) {KKT/Schur};
\node[b, right=of d] (e) {Baumgarte};
\node[b, below=6mm of a] (f) {闭链};
\node[b, right=of f] (g) {接触/CWC};
\node[b, right=of g] (h) {Gauss/UK/IFT};
\node[b, right=of h] (i) {可微接触};
\node[b, below=6mm of f] (j) {$\partial$RNEA $O(Nd)$};
\node[b, right=of j] (k) {$\partial$ABA/AZA};
\node[b, right=of k] (l) {二阶/辨识};
\node[b, right=of l] (m) {DDP/MPC};
\draw[ln] (a)--(c);\draw[ln] (c)--(d);\draw[ln] (d)--(e);
\draw[ln] (e.south)-- ++(0,-2.5mm) -| (f.north);
\draw[ln] (f)--(g);\draw[ln] (g)--(h);\draw[ln] (h)--(i);
\draw[ln] (i.south)-- ++(0,-2.5mm) -| (j.north);
\draw[ln] (j)--(k);\draw[ln] (k)--(l);\draw[ln] (l)--(m);
\end{tikzpicture}
\end{center}

\paragraph{启下。}
本章是「刚体动力学」子部（\cref{ch:spatial-dynamics} 空间向量地基 $\to$ \cref{ch:geometric-mechanics} 几何力学 $\to$ 本章约束）的收口，向外伸出多条钩子：辛/几何力学的深化（Baumgarte 破辛、SHAKE/RATTLE 保辛）回到 \cref{ch:geometric-mechanics}；接触的求解器与接触力分配主体在 \cref{ch:contact-mechanics}；操作空间动力学的 $\boldsymbol\Lambda=\mathbf G^{-1}$ 与阻抗在 \cref{ch:operational-space}；WBC/TSID-HQP 的库实现在 \cref{ch:force-wbc}（其理论根即本章 \cref{sec:cd-gauss} 的「WBC$=$Gauss 多任务推广」）；解析导数装配的 DDP/iLQR 在 \cref{ch:ddp-ilqr}，proximal 约束动力学喂实时 MPC 在 \cref{ch:mpc-advanced}；接触隐式 MPC 与可微仿真也是规划部路径积分/MPPI 的前向仿真基础。\cref{ch:control} 机械臂部的并联闭链将特化本章 \cref{sec:cd-loop}（闭链 KKT）/\cref{sec:cd-kkt}（Delassus）/\cref{sec:cd-drnea}--\cref{sec:cd-daba}（解析导数）到 Delta/Stewart 与 $7$ 自由度操作空间约束；腿足部则引用本章浮基方程、CWC/ZMP/CP 与接触 LCP 框架。沿这些接口走下去，约束动力学就从「会写方程」变成「能在真实机器人上以千赫兹求解、以微秒求导、以梯度优化」的活工具。

% ════════════════════════════════════════════════════════════════
\section*{常见误解汇总}
\addcontentsline{toc}{section}{常见误解汇总}
\begin{itemize}
  \item \textbf{「nonholonomic = 到不了某些位置」}：错。Chow--Rashevsky 保证 bracket-generating 系统可达全空间，约束限的是瞬时速度方向（\cref{sec:cd-classify}）。
  \item \textbf{「高阶积分器能消除 drift-off」}：错。对 index-3 DAE 任何标准 ODE 积分器都漂移，只是速率不同（\cref{thm:cd-drift}）。
  \item \textbf{「约束力是 $\mathbf J\boldsymbol\lambda$」}：错。永远是 $\mathbf J^\top\boldsymbol\lambda$，缺转置力方向错（\cref{sec:cd-kkt}）。
  \item \textbf{「Delassus 算子是 $\boldsymbol\Lambda=(\mathbf J\mathbf M^{-1}\mathbf J^\top)^{-1}$」}：不精确。Delassus 矩阵是 $\mathbf G=\mathbf J\mathbf M^{-1}\mathbf J^\top$ 本身，其逆才是 Khatib OSIM $\boldsymbol\Lambda$（\cref{thm:cd-schur}）。
  \item \textbf{「Baumgarte $\alpha$ 越大越好」}：错。$\alpha h>1$ 时显式积分发散（\cref{sec:cd-baumgarte}）。
  \item \textbf{「4 边摩擦锥够用」}：对角方向比圆锥少约 $29\%$，大 $\mu_f$ 须 8-polygon 或 SOCP（\cref{sec:cd-contact}）。
  \item \textbf{「CWC yaw 界是 $|\tau_z|\le\mu_f(X+Y)f_z$」}：那只是最松界，紧界与切向耦合（\cref{eq:cd-cwc-yaw}）。
  \item \textbf{「$\partial\boldsymbol\tau/\partial\mathbf q$ 是下三角稀疏」}：错。串联链通常稠密，$O(Nd)$ 是算法复杂度而非结果稀疏（\cref{sec:cd-drnea}）。
\end{itemize}

\section*{小结}
\addcontentsline{toc}{section}{小结}

\begin{table}[H]\centering\small
\caption{本章核心结论速查}
\label{tab:cd-summary}
\begin{tabular}{p{0.20\linewidth}p{0.46\linewidth}p{0.22\linewidth}}
\toprule
知识点 & 核心结果 & 出处 \\
\midrule
约束分类 & Pfaffian $+$ Frobenius/Signorini/Chow；对合 $\Leftrightarrow$ 可积证、$\omega\wedge\mathrm d\omega$ 判据 & \cref{thm:cd-frobenius,lem:cd-frobenius-1form} \\
DAE 指数 & 位置 $3$/速度 $2$/加速度 $1$；drift-off $O(t^2)$ & \cref{tab:cd-ladder,thm:cd-drift} \\
KKT 系统 & 鞍点 $+$ Schur，$\mathbf G=\mathbf J\mathbf M^{-1}\mathbf J^\top$、$\boldsymbol\Lambda=\mathbf G^{-1}$ & \cref{eq:cd-kkt,thm:cd-schur} \\
Baumgarte & $\ddot e+2\alpha\dot e+\beta^2 e=0$，临界 $\alpha=\beta$、$\alpha h\le1$ & \cref{eq:cd-baumgarte-error} \\
闭链 & cut-joint $+$ Schur 补，$O(N)+O(m^3)$ & \cref{eq:cd-loop-schur} \\
接触 & Signorini $+$ Coulomb $+$ LCP/Stewart--Trinkle/Moreau & \cref{eq:cd-signorini,thm:cd-stewart} \\
Gauss/UK & 最小约束原理 $\Rightarrow$ UK 显式解；IFT-on-KKT & \cref{thm:cd-gauss,eq:cd-uk,eq:cd-ift-kkt} \\
$\partial$RNEA & placement 恒等式 $\Rightarrow$ $O(Nd)$ 解析递推 & \cref{thm:cd-placement,alg:cd-drnea} \\
$\partial$ABA & ID--FD 互逆 $\Rightarrow\partial\ddot{\mathbf q}/\partial\mathbf q=-\mathbf M^{-1}\partial\boldsymbol\tau/\partial\mathbf q$；AZA & \cref{eq:cd-idfd,eq:cd-aza} \\
DDP Jacobian & $\mathbf A=[\mathbf 0,\mathbf I;\,\partial\ddot{\mathbf q}/\partial\mathbf q,\partial\ddot{\mathbf q}/\partial\mathbf v]$，$\mathbf B=[\mathbf 0;\mathbf M^{-1}]$ & \cref{eq:cd-ddp-jac} \\
\bottomrule
\end{tabular}
\end{table}

\section*{延伸阅读}
\addcontentsline{toc}{section}{延伸阅读}
\begin{itemize}
  \item \textbf{约束/闭链动力学}：Featherstone, \emph{Rigid Body Dynamics Algorithms}\cite{featherstone2008rigid}（第 8 章闭链）；Sathya \& Carpentier 的 constrainedABA\cite{sathya2024constrained}（$O(N)$ 约束前向动力学）。
  \item \textbf{DAE 与稳定化}：Baumgarte 1972\cite{baumgarte1972stabilization}（原典）；Petzold 1982\cite{petzold1982dae} 与 Hairer--Wanner\cite{hairer1996solving}（DAE 指数/drift-off）；GGL\cite{gear1985automatic}、SHAKE/RATTLE\cite{ryckaert1977numerical,andersen1983rattle}。
  \item \textbf{接触与非光滑力学}：Painlevé 1895\cite{painleve1895lois}；Stewart--Trinkle 1996\cite{stewart1996implicit}、Moreau 1988\cite{moreau1988unilateral}、Anitescu 2006\cite{anitescu2006optimization}；Brogliato, \emph{Nonsmooth Mechanics}\cite{brogliato2016nonsmooth}（全书）；Castro 等 SAP\cite{castro2021unconstrained}；Le Lidec 等接触模型比较\cite{lelidec2024contact}。
  \item \textbf{稳定性判据}：ZMP\cite{vukobratovic1972stability}、CWC\cite{caron2015stability}、Capture Point\cite{pratt2006capture}、质心动力学\cite{orin2013centroidal}。
  \item \textbf{统一原理与可微}：Udwadia--Kalaba\cite{udwadia1992new}（含 Gauss 原理）；Khatib 操作空间\cite{khatib1987unified}；proximal 约束动力学\cite{carpentier2021proximal}；可微物理 Dojo\cite{howell2022dojo}、Nimble\cite{werling2021fast}、随机平滑\cite{suh2022differentiable}、MuJoCo\cite{todorov2014convex}；约束/欠驱动操作空间 Mistry--Righetti\cite{mistry2012operational}。
  \item \textbf{解析微分}：Carpentier--Mansard\cite{carpentier2018analytical}（一阶 $\partial$RNEA）与 JSIM 解析逆\cite{carpentier2018inverse}；Singh--Russell--Wensing 一阶（AZA）\cite{singh2021efficient} 与二阶\cite{singh2022second}；惯性辨识 Khalil--Dombre\cite{khalil2002modeling}。
\end{itemize}



